\documentclass[11pt,a4paper]{article}

% Encoding and fonts
\usepackage[utf8]{inputenc}
\usepackage[T1]{fontenc}
\usepackage{lmodern}

% Page geometry
\usepackage[margin=2.5cm]{geometry}

% Math
\usepackage{amsmath,amssymb,amsfonts}

% Tables
\usepackage{longtable}
\usepackage{booktabs}
\usepackage{array}

% Graphics and colors
\usepackage{graphicx}
\usepackage{adjustbox}
\usepackage[dvipsnames]{xcolor}

% Hyperlinks
\usepackage[colorlinks=true,linkcolor=NavyBlue,urlcolor=NavyBlue,citecolor=NavyBlue]{hyperref}

% Code / verbatim
\usepackage{fancyvrb}
\usepackage{listings}
\lstset{
  basicstyle=\ttfamily\small,
  breaklines=true,
  frame=single,
  backgroundcolor=\color{gray!5},
  xleftmargin=1em,
  framexleftmargin=0.5em
}

% Spacing
\usepackage{setspace}
\onehalfspacing
\usepackage{parskip}

% Headers
\usepackage{fancyhdr}
\pagestyle{fancy}
\fancyhf{}
\fancyhead[L]{\small\nouppercase{\leftmark}}
\fancyhead[R]{\small\thepage}
\renewcommand{\headrulewidth}{0.4pt}

% Title
\title{\LARGE\textbf{Research Programme Overview}\\[0.5em]
       \Large Non-Bank Financial Intermediation and Systemic Risk}
\author{Sergio Sola}
\date{February 2026}

\begin{document}
\maketitle
\thispagestyle{empty}

\begin{abstract}
\noindent This document provides a detailed research cookbook for a four-project
programme studying how non-bank financial intermediation (NBFI) creates, amplifies,
and transmits systemic risk. For each project, every analytical module specifies:
(1)~variables---exact names, definitions, sources, frequency, and level of aggregation;
(2)~model specification---full equation with every term defined;
(3)~estimation---method, software, key parameters;
(4)~output---what you get and how to interpret it; and
(5)~module linkages---how each module's output feeds into the next.
The unifying theme is that NBFI risks are fundamentally nonlinear: largely invisible
in normal times but activating powerfully during stress.
\end{abstract}

\tableofcontents
\newpage


\section{Research Programme Overview: Non-Bank Financial Intermediation and Systemic Risk}

\textbf{Author:} Sergio Sola | February 2026


\bigskip\noindent\rule{\textwidth}{0.4pt}\bigskip

This document is a \textbf{detailed research cookbook} for the four-project programme. For each project, every analytical module specifies:

\begin{enumerate}
\item \textbf{Variables} --- exact names, definitions, sources, frequency, and level of aggregation
\item \textbf{Model specification} --- full equation with every term defined
\item \textbf{Estimation} --- method, software, key parameters (lags, quantiles, etc.)
\item \textbf{Output} --- what you get and how to interpret it
\item \textbf{Module linkages} --- how each module's output feeds into the next

\end{enumerate}
\textbf{Unifying theme:} NBFI risks are fundamentally nonlinear --- largely invisible in normal
times but activating powerfully during stress. Standard mean-based methods miss these
dynamics. All four projects employ quantile-specific methods to capture tail amplification
and asymmetric contagion.


{\small
\begin{longtable}{|p{0.22\textwidth}|p{0.33\textwidth}|p{0.33\textwidth}|}
\hline
\textbf{Project} & \textbf{Title} & \textbf{Focus} \\
\hline
\endhead
P1 & The Shadow Leverage Map & Hidden NBFI leverage from repo \& derivatives \\
\hline
P2 & Mapping the Channels & NBFI amplification of monetary policy transmission \\
\hline
P3 & Contagion Across Borders & Cross-border NBFI spillovers via the GFC \\
\hline
P4 & FX Hedging as a Contagion Channel & FX hedging costs as a transmission mechanism \\
\hline
\end{longtable}
}


\bigskip\noindent\rule{\textwidth}{0.4pt}\bigskip


\subsection{How to Read This Document}

Each project is organised into \textbf{numbered modules}. Every module follows this template:


\begin{quote}
\itshape
\textbf{Purpose} --- what economic question this module answers \\

\textbf{Unit of observation} --- e.g., ''sub-sector \$s\$ $\times$ quarter \$t\$'' \\

\textbf{Input variables} --- table listing every variable entering the model, its definition, source, frequency, and level of aggregation \\

\textbf{Model specification} --- the full equation with all terms defined \\

\textbf{Estimation details} --- how the model is estimated (OLS, quantile regression, MLE, simulation, etc.), plus key parameters \\

\textbf{Output} --- what the function returns and how to interpret the results \\

\textbf{$\rightarrow$ Link to next module} --- what this module's output feeds into \\
\end{quote}

At the end of each project there is a \textbf{Module Linkage Map} showing the pipeline flow.


\bigskip\noindent\rule{\textwidth}{0.4pt}\bigskip


\section{Project 1: The Shadow Leverage Map}

\textbf{Subtitle:} Tracing Hidden Non-Bank Exposures Through Repo and Derivatives Data


\subsection{Motivation}

NBFIs now account for nearly half of global financial assets, yet their leverage
remains largely invisible to regulators. Unlike banks, NBFIs are not subject to
standardised capital or leverage ratio disclosure, creating a systemic blind spot.
This project develops a \textit{Shadow Leverage Map} --- a prototype analytical framework
that reverse-engineers the leverage embedded in NBFI balance sheets using publicly
available market data from repo markets (OFR, ECB MMSR), derivatives reporting
(DTCC), and FSB monitoring data.


\subsection{Research Questions}

\textbf{RQ1.} Can NBFI leverage be reverse-engineered from observable market data
(repo volumes, derivatives notional, equity proxies)?

\textbf{RQ2.} How are banks and NBFIs interconnected through repo and derivatives
exposures, and how do these interconnections evolve over time and during stress?

\textbf{RQ3.} Can a spectral measure of network structure --- the Hidden Leverage Index
(HLI), based on the largest eigenvalue of the exposure-weighted adjacency matrix
--- serve as a predictor of systemic events?

\textbf{RQ4.} Does hidden leverage amplify tail risk during stress? Is the bank-NBFI
system more interconnected in the tails (quantile connectedness), and does high
hidden leverage predict worse left-tail growth (Growth-at-Risk)?


\subsection{Empirical Strategy}

Project 1 has \textbf{four modules} that build on each other in sequence:


\begin{adjustbox}{max width=\textwidth}
\begin{minipage}{1.5\textwidth}
\begin{verbatim}
Module 1 (Implied Leverage)  --→  Module 2 (Bipartite Network)  --→  Module 3 (HLI Spectral Index)  --→  Module 4 (Tail Risk)
         ↓                                  ↓                                  ↓                                ↓
   Leverage ratios              Exposure-weighted               Scalar HLI_t time              HLI predicts left-tail
   per NBFI sector              adjacency matrix W_t            series (leading              GDP growth (GaR) and
   per quarter                  evolving quarterly              indicator)                    quantile connectedness
\end{verbatim}
\end{minipage}
\end{adjustbox}


\bigskip\noindent\rule{\textwidth}{0.4pt}\bigskip


\subsubsection{Module 1: Implied Leverage Estimation}

\textbf{Purpose}: Reverse-engineer the economic leverage embedded in NBFI balance sheets --- leverage that regulatory ratios do not capture --- using observable market data.

\textbf{Unit of observation}: NBFI sub-sector \$s\$ $\times$ quarter \$t\$.

\textbf{Input variables}:


{\small
\begin{longtable}{|p{0.15\textwidth}|p{0.15\textwidth}|p{0.15\textwidth}|p{0.15\textwidth}|p{0.15\textwidth}|p{0.15\textwidth}|}
\hline
\textbf{Variable} & \textbf{Symbol} & \textbf{Definition} & \textbf{Source} & \textbf{Frequency} & \textbf{Level} \\
\hline
\endhead
Equity proxy & $E_{s,t}$ & Total net asset value (equity) of NBFI sub-sector \$s\$ & FSB Global Monitoring Report; Fed Z.1 Flow of Funds & Annual (FSB); Quarterly (Z.1) & Sub-sector aggregate \\
\hline
Repo borrowing & $R_{s,t}$ & Total repo borrowing by sub-sector \$s\$ & OFR US Repo (SOFR data); ECB MMSR & Daily (aggregated to quarterly) & Sub-sector aggregate \\
\hline
Derivatives notional & $D_{s,t}$ & Gross derivatives notional outstanding for sub-sector \$s\$ & DTCC Swap Data Repository & Weekly (aggregated to quarterly) & Sub-sector aggregate \\
\hline
Delta-equivalence factor & $\delta$ & Scaling factor converting gross notional to economic exposure & Calibrated from BIS Triennial Survey & Fixed parameter & Global \\
\hline
\end{longtable}
}

\textbf{Sub-sectors} $s \in \{$hedge funds, MMFs, bond mutual funds, equity mutual funds, ETFs, broker-dealers, insurance companies, pension funds, finance companies$\}$.

\textbf{Model specification}:

\[\text{Implied Leverage}_{s,t} = \frac{E_{s,t} + R_{s,t} + \delta \cdot D_{s,t}}{E_{s,t}}\]

\textbf{Worked example}: A hedge fund sector with $E = \$10\text{B}\$ equity, $R = \$40\text{B}\$ repo borrowing, and $D = \$500\text{B}\$ IRS notional (with $\delta = 0.05$) has implied leverage of $(10 + 40 + 0.05 \times 500)/10 = 7.5\times$.

\textbf{Estimation details}:
\begin{itemize}
\item This is a \textit{constructed measure}, not an econometric model --- no estimation required.
\item The key calibration choice is $\delta$. Baseline: $\delta = 0.05$ (from BIS survey data on delta-equivalence of IRS portfolios). Robustness: $\delta \in [0.02, 0.10]$.
\item For sub-sectors where direct repo data is unavailable, repo borrowing is proxied by total short-term funding minus equity, using Flow of Funds (Z.1) liability breakdowns.

\end{itemize}
\textbf{Output}: A panel of implied leverage ratios: $\text{Lev}_{s,t}$ for each sub-sector \$s\$ and each quarter \$t\$ (2013Q1--2024Q4). This is a single number per sub-sector per quarter --- the ''true'' economic leverage.

\textbf{$\rightarrow$ Link to Module 2}: The implied leverage ratios serve as \textit{edge weights} in the bipartite network. Higher leverage for sub-sector $s$ means the edges from banks to $s$ carry more systemic weight (because a given \$ exposure to a highly-leveraged NBFI poses more fire-sale risk).


\bigskip\noindent\rule{\textwidth}{0.4pt}\bigskip


\subsubsection{Module 2: Bipartite Bank-NBFI Network}

\textbf{Purpose}: Map the bilateral exposure network between G-SIB banks and NBFI sub-sectors, weighted by both exposure amounts and the implied leverage from Module 1.


\paragraph{Geographic Dimension of Exposures}

The economically correct unit of observation is:

\[\text{Exp}_{n,i,j,m,t} = \text{exposure of bank } i \text{ in country } n \text{ to NBFI sub-sector } j \text{ in country } m \text{ at time } t\]

However, bilateral bank-by-NBFI-by-country data is generally unavailable. The realistic data structure is \textbf{semi-bilateral}: for each NBFI sub-sector \$j\$ in country \$n\$, we observe total cross-border exposure to the rest of the world, $\text{Exp}_{j,n,t}^{\text{RoW}}$, but not the specific \$(n,m)\$ counterparty breakdown at the sub-sector level. This means:

\begin{itemize}
\item \textbf{Domestic exposures}: Where bilateral data exists (e.g., US via OFR, EU via ECB MMSR), we can construct full bank-by-NBFI matrices within a jurisdiction.
\item \textbf{Cross-border exposures}: Available only as (sub-sector \$j\$, country \$n\$) $\rightarrow$ rest-of-world aggregates (from BIS Locational Banking Statistics, which report bank claims on/liabilities to non-bank financial institutions by country, but not by NBFI sub-type across borders).

\end{itemize}
In practice, $\mathcal{W}_t$ is constructed as a \textbf{two-tier structure}:
\begin{enumerate}
\item \textbf{Within-jurisdiction}: Full \$(i,j)\$ bilateral matrix where OFR/ECB data permits.
\item \textbf{Cross-border}: $(j,n) \to \text{RoW}$ aggregate exposures, allocated proportionally using BIS country-level banking statistics as weights.

\end{enumerate}
This is an honest limitation. We can still compute system-wide amplification ($\lambda_1$), but the cross-border contagion paths are identified at the country level rather than the bank-NBFI pair level. The geographic decomposition connects directly to Project 3 (cross-border contagion).


\paragraph{Input variables}


{\small
\begin{longtable}{|p{0.15\textwidth}|p{0.15\textwidth}|p{0.15\textwidth}|p{0.15\textwidth}|p{0.15\textwidth}|p{0.15\textwidth}|}
\hline
\textbf{Variable} & \textbf{Symbol} & \textbf{Definition} & \textbf{Source} & \textbf{Frequency} & \textbf{Level} \\
\hline
\endhead
Bilateral exposure (domestic) & $\text{Exp}_{ij,t}$ & Bank \$i\$'s total exposure to NBFI sub-sector \$j\$ (within jurisdiction) & OFR bilateral repo data (US); ECB MMSR (EU) & Quarterly & Bank $\times$ NBFI sub-sector \\
\hline
Cross-border exposure & $\text{Exp}_{j,n,t}^{\text{RoW}}$ & Total exposure of NBFI sub-sector \$j\$ in country \$n\$ to rest-of-world & BIS Locational Banking Stats & Quarterly & Sub-sector $\times$ country \\
\hline
Implied leverage & $\text{Lev}_{j,t}$ & From Module 1 & Module 1 output & Quarterly & NBFI sub-sector \\
\hline
Institution metadata & --- & Bank names, NBFI sub-sector classification, country & FSB G-SIB list; own classification & Static & Institution \\
\hline
\end{longtable}
}


\paragraph{Network construction}

The network is a \textbf{directed weighted bipartite graph} $G_t = (V_B \cup V_N, E_t)$ where $V_B$ = G-SIB banks, $V_N$ = NBFI sub-sectors, and $E_t$ = directed edges representing credit and funding relationships.


\paragraph{Construction of the Adjacency Matrix $\mathcal{W}_t$}

The \textbf{exposure-weighted adjacency matrix} is:

\[W_{ij,t}^{\text{credit}} = \frac{\text{Exp}_{ij,t} \times \text{Lev}_{j,t}}{\sum_k \text{Exp}_{ik,t} \times \text{Lev}_{k,t}}\]

\textbf{Why multiply by leverage?} The logic is a \textbf{fire-sale argument}. Consider bank \$i$ that has lent \$100M to NBFI sub-sector \$j\$ via repo. If \$j\$ is leveraged 2$\times$, a 10\% asset price decline wipes out 20\% of its equity --- painful but survivable. If \$j\$ is leveraged 10$\times$, the same 10\% decline wipes out 100\% of equity --- triggering forced liquidation, fire sales, and potential default. The loss that cascades back to bank \$i\$ is therefore proportional not just to the raw exposure $\text{Exp}_{ij,t}$ but to the \textit{leverage-adjusted exposure} $\text{Exp}_{ij,t} \times \text{Lev}_{j,t}$. Leverage acts as an amplifier: the same dollar of exposure to a highly-leveraged counterparty poses far more systemic risk than to a low-leverage one.

The row-normalisation (dividing by $\sum_k$) ensures each row sums to 1, giving the \textit{share} of bank $i$'s leverage-weighted risk that comes from each NBFI sub-sector.


\paragraph{Two Directions of Risk: Credit vs. Funding}

The matrix $\mathcal{W}_t^{\text{credit}}$ (dimension $n_B \times n_N$) captures only one direction --- \textbf{credit/fire-sale risk} (bank's loss from NBFI distress). There is a second, equally important direction --- \textbf{funding risk} (NBFI's loss from bank funding withdrawal). The funding matrix $\mathcal{W}_t^{\text{funding}}$ has dimension $n_N \times n_B$, with $(j,i)$-th entry:

\[W_{ji,t}^{\text{funding}} = \frac{\text{Exp}_{ji,t} \times \text{Lev}_{i,t}}{\sum_k \text{Exp}_{jk,t} \times \text{Lev}_{k,t}}\]

This captures: NBFI sub-sector \$j\$ depends on bank \$i\$ for repo funding or credit lines. If bank \$i\$ is highly leveraged and its own VaR constraint binds, it pulls back funding (raises haircuts, calls margin, refuses to roll repo). The severity depends on how leveraged bank \$i\$ is --- more leveraged banks are more likely to pull funding when stressed.


\paragraph{The Combined System Matrix}

For the spectral analysis (Module 3), we need a \textbf{single} matrix that captures the feedback loop between credit and funding risk: bank pulls funding $\rightarrow$ NBFI fire-sells $\rightarrow$ bank takes credit loss $\rightarrow$ bank pulls more funding. Two separate matrices would miss this circularity.

We combine both into a single $(n_B + n_N) \times (n_B + n_N)$ block matrix:

\[\mathcal{W}_t = \begin{pmatrix} 0 & W_t^{\text{credit}} \\ W_t^{\text{funding}} & 0 \end{pmatrix}\]

where:
\begin{itemize}
\item The upper-right block $\mathcal{W}_t^{\text{credit}}$ ($n_B \times n_N$): bank \$i\$'s credit risk exposure to NBFI \$j\$
\item The lower-left block $\mathcal{W}_t^{\text{funding}}$ ($n_N \times n_B$): NBFI \$j\$'s funding risk from bank \$i\$
\item The diagonal blocks are zero (within-bank and within-NBFI linkages are not modelled here)

\end{itemize}
The spectral radius $\lambda_1(\mathcal{W}_t)$ of this combined matrix captures the \textbf{total amplification potential} from both channels feeding into each other. A high $\lambda_1$ means the credit-funding feedback loop is tightly wound --- NBFI fire sales hurt banks, which pull funding, which triggers more fire sales.


\paragraph{NBFI Counterparty Decomposition}

A critical refinement: the systemic risk implications of NBFI activity depend fundamentally on \textbf{who the NBFI transacts with}. An NBFI sub-sector's balance sheet should be decomposed by counterparty:


{\small
\begin{longtable}{|p{0.22\textwidth}|p{0.33\textwidth}|p{0.33\textwidth}|}
\hline
\textbf{Domestic banks} & \textbf{Foreign counterparties} & \textbf{Domestic private sector} \\
\hline
\endhead
\textbf{NBFI assets (lends to)} & Repo lending, deposits, bank bonds $\rightarrow$ NBFI is \textit{funding source} for banks & Cross-border portfolio holdings, FX swaps $\rightarrow$ NBFI transmits shocks \textit{outward} \\
\hline
\textbf{NBFI liabilities (borrows from)} & Repo borrowing, credit lines, prime brokerage $\rightarrow$ Bank exposed to NBFI \textit{credit risk} & Cross-border funding, FX swap borrowing $\rightarrow$ NBFI \textit{imports} foreign funding shocks \\
\hline
\end{longtable}
}

Each cell creates a different contagion channel with different speed, amplification, and policy implications. A country where NBFIs primarily borrow from domestic banks and lend to the private sector (credit intermediation model) has a completely different risk profile from one where NBFIs borrow from foreigners and trade derivatives with domestic banks (market intermediation model).

\textbf{Data source}: ''From-whom-to-whom'' financial accounts --- Fed Z.1 (US), ECB Financial Accounts (euro area), OECD institutional sector accounts (broader coverage).

\textbf{Implication for the adjacency matrix}: Where from-whom-to-whom data is available, $\mathcal{W}_t$ can be disaggregated by counterparty type, enabling separate spectral analyses for credit-intermediation vs. market-intermediation risk channels.


\paragraph{Estimation details}

\begin{itemize}
\item No econometric estimation --- this is a data construction step.
\item Network is built using \texttt{networkx} (Python). The function \texttt{build\_exposure\_network()} creates a \texttt{DiGraph} from the exposure panel.
\item Centrality measures computed: PageRank, eigenvector centrality, betweenness centrality, in/out-strength (weighted degree).
\item The network evolves quarterly from 2015Q1 to 2024Q4 (limited by BIS sectoral breakdown availability; extended to 2013 where OFR bilateral data permits).

\end{itemize}

\paragraph{Output}

\begin{itemize}
\item A time series of adjacency matrices $\mathcal{W}_t^{\text{credit}}$ ($n_B \times n_N$) and $\mathcal{W}_t^{\text{funding}}$ ($n_N \times n_B$), one pair per quarter.
\item Combined system matrix $\mathcal{W}_t$ ($(n_B + n_N) \times (n_B + n_N)$ block matrix) for spectral analysis.
\item Centrality rankings for each node at each date.
\item Network summary statistics over time: density, total exposure, HHI (concentration), reciprocity.
\item Decomposition of cross-border vs. domestic exposure shares over time.

\end{itemize}
\textbf{$\rightarrow$ Link to Module 3}: The combined adjacency matrix $\mathcal{W}_t$ is the direct input to the spectral analysis. The HLI is the largest eigenvalue of $\mathcal{W}_t$.


\bigskip\noindent\rule{\textwidth}{0.4pt}\bigskip


\subsubsection{Module 3: Hidden Leverage Index (Spectral)}

\textbf{Purpose}: Collapse the bipartite network into a single scalar measure of systemic amplification potential --- the Hidden Leverage Index (HLI) --- using the spectral radius of the exposure matrix.

\textbf{Unit of observation}: Aggregate system-level measure $\times$ quarter \$t\$.

\textbf{Input variables}:


{\small
\begin{longtable}{|p{0.15\textwidth}|p{0.15\textwidth}|p{0.15\textwidth}|p{0.15\textwidth}|p{0.15\textwidth}|p{0.15\textwidth}|}
\hline
\textbf{Variable} & \textbf{Symbol} & \textbf{Definition} & \textbf{Source} & \textbf{Frequency} & \textbf{Level} \\
\hline
\endhead
Combined adjacency matrix & $\mathcal{W}_t$ & Combined credit + funding block matrix from Module 2 & Module 2 output & Quarterly & System-level \\
\hline
\end{longtable}
}

\textbf{Model specification}:

\[\text{HLI}_t = \lambda_1(\mathcal{W}_t)\]

where $\lambda_1(\mathcal{W}_t)$ is the \textbf{largest eigenvalue} (spectral radius) of the combined system matrix $\mathcal{W}_t$.

\textbf{Economic intuition --- network amplification, not PCA} (Acemoglu, Ozdaglar \& Tahbaz-Salehi 2015): The spectral radius is \textit{not} a statistical summary like PCA's first component --- it is a \textbf{structural measure of the network's amplification capacity}. The logic is a cascade argument:

Consider a shock $\varepsilon$ hitting one node. Through the network, losses propagate: node $j$'s distress hits bank $i$ with weight $\mathcal{W}_{ij}$, then $i$'s distress hits others with weight $\mathcal{W}_{ki}$, and so on. After $n$ rounds, total losses are:

\[\text{Total loss} \propto \varepsilon + \mathcal{W}\varepsilon + \mathcal{W}^2\varepsilon + \mathcal{W}^3\varepsilon + \ldots = (I - \mathcal{W})^{-1}\varepsilon\]

This geometric series converges if and only if $\lambda_1(\mathcal{W}) < 1$. The \textbf{rate of amplification} is governed by $\lambda_1$:
\begin{itemize}
\item $\lambda_1$ close to 0: shocks dissipate quickly. The network is ''loose.''
\item $\lambda_1$ close to 1: shocks amplify through many rounds before dying out. The network is ''tightly wound'' --- approaching the instability threshold.

\end{itemize}
The \textbf{eigenvector} associated with $\lambda_1$ also matters --- it identifies \textbf{which nodes are most central to the amplification}. If the leading eigenvector loads heavily on hedge funds, that is where systemic risk concentrates.

\textbf{Estimation details}:
\begin{itemize}
\item Computed via \texttt{numpy.linalg.eigvals} on $\mathcal{W}_t$ at each quarter.
\item No estimation --- this is a deterministic transformation of $\mathcal{W}_t$.
\item Robustness: also compute the second eigenvalue $\lambda_2$ (captures secondary amplification channels) and the spectral gap $\lambda_1 - \lambda_2$ (captures how ''dominant'' the main amplification channel is).
\item Separate spectral analysis on $W^{\text{credit}}$ vs. $W^{\text{funding}}$ to identify which direction of risk dominates amplification.

\end{itemize}
\textbf{Output}: A quarterly time series $\{\text{HLI}_t\}_{t=2013Q1}^{2024Q4}$ --- a single number per quarter measuring systemic amplification potential.

\textbf{Validation}: The HLI should spike \textit{before} known stress episodes:
\begin{itemize}
\item March 2020 (COVID dash-for-cash)
\item September 2022 (UK gilt/LDI crisis)
\item March 2023 (SVB failure)

\end{itemize}
If it does, this confirms its value as a leading indicator.

\textbf{$\rightarrow$ Link to Module 4}: The HLI time series enters Module 4 as a predictor in tail risk regressions (quantile connectedness and Growth-at-Risk).


\bigskip\noindent\rule{\textwidth}{0.4pt}\bigskip


\subsubsection{Module 4: Tail Risk Econometrics}

Module 4 has \textbf{three sub-analyses}, all using the HLI from Module 3 as a key variable.


\bigskip\noindent\rule{\textwidth}{0.4pt}\bigskip


\paragraph{Module 4a: Quantile Connectedness (Bank-NBFI System)}

\textbf{Purpose}: Test whether the bank-NBFI system is more interconnected in the tails (stress) than at the median (normal times).

\textbf{Unit of observation}: System of \$k\$ return series $\times$ week \$t\$.

\textbf{Input variables}:


{\small
\begin{longtable}{|p{0.15\textwidth}|p{0.15\textwidth}|p{0.15\textwidth}|p{0.15\textwidth}|p{0.15\textwidth}|p{0.15\textwidth}|}
\hline
\textbf{Variable} & \textbf{Symbol} & \textbf{Definition} & \textbf{Source} & \textbf{Frequency} & \textbf{Level} \\
\hline
\endhead
Bank sector returns & $r_{i,t}^B$ & Weekly log-returns for G-SIB banks (or bank sector index) & Bloomberg / Refinitiv & Weekly & Institution or sector \\
\hline
NBFI sector returns & $r_{j,t}^N$ & Weekly log-returns for NBFI sub-sector indices (HF index, MMF index, insurance index, etc.) & Bloomberg (HFRX, etc.) & Weekly & Sub-sector index \\
\hline
\end{longtable}
}

\textbf{System}: A vector $Y_t = (r_1^B, \ldots, r_m^B, r_1^N, \ldots, r_n^N)'$ of \$k = m + n\$ return series (e.g., 5 bank sector returns + 5 NBFI sub-sector returns = 10-variable system).

\textbf{Model specification} --- Quantile VAR at quantile $\tau$:

\[Q_{\tau}(y_{i,t} \mid Y_{t-1}, \ldots, Y_{t-p}) = c_i(\tau) + \sum_{j=1}^{k} \sum_{l=1}^{p} \beta_{ij}^{(l)}(\tau) \, y_{j,t-l}\]

This is a standard VAR except that \textit{every coefficient depends on the quantile $\tau$}. At $\tau = 0.05$, you are modelling the left tail; at $\tau = 0.50$, the median; at $\tau = 0.95$, the right tail.

\textbf{Estimation details}:
\begin{itemize}
\item Estimated \textbf{equation by equation} using \texttt{statsmodels.QuantReg} (linear quantile regression, Koenker \& Bassett 1978).
\item For each equation \$i\$: the dependent variable is $y_{i,t}$; the regressors are \$p\$ lags of all \$k\$ variables plus an intercept.
\item Lags: \$p = 4\$ (quarterly data) or \$p = 2\$ (weekly data). Selected via BIC at each quantile.
\item Quantiles estimated: $\tau \in \{0.05, 0.25, 0.50, 0.75, 0.95\}$.
\item From the estimated QVAR, compute the \textbf{Generalised Forecast Error Variance Decomposition (GFEVD)} at horizon \$h = 10\$:
\item Uses the companion matrix to compute MA representation.
\item GFEVD entry $\theta_{ij}(\tau)$ = share of variable \$i\$'s \$h\$-step forecast error variance attributable to shocks in variable \$j\$, at quantile $\tau$.
\item Normalised so rows sum to 100\%.

\end{itemize}
\textbf{Connectedness measures} (Ando, Greenwood-Nimmo \& Shin 2022):
\begin{itemize}
\item \textbf{Total connectedness}: $C(\tau) = \frac{1}{k} \sum_{i \neq j} \theta_{ij}(\tau) \times 100\%$
\item \textbf{To-others}: $\text{TO}_i(\tau) = \sum_{j \neq i} \theta_{ji}(\tau)$ --- how much \$i\$ transmits to the system
\item \textbf{From-others}: $\text{FROM}_i(\tau) = \sum_{j \neq i} \theta_{ij}(\tau)$ --- how much \$i\$ receives from the system
\item \textbf{Net}: $\text{NET}_i(\tau) = \text{TO}_i - \text{FROM}_i$ --- net transmitter (\$>0\$) or net receiver (\$<0\$)

\end{itemize}
\textbf{Output}:
\begin{enumerate}
\item A $k \times k$ GFEVD matrix $\Theta(\tau)$ at each quantile --- shows pairwise spillover direction and magnitude.
\item Total connectedness at each quantile: \$C(0.05)\$, \$C(0.50)\$, \$C(0.95)\$.
\item Directional spillovers: which sectors are net transmitters vs. receivers of risk, and how this changes between median and tails.

\end{enumerate}
\textbf{Key test}: Is $C(0.05) \gg C(0.50)$? If total connectedness at the 5th percentile is much higher than at the median, the system is asymmetrically interconnected --- more contagious during stress.

\textbf{Rolling version}: Estimated on a rolling window (60 weeks), stepping 1 week at a time. Produces a time series of $C_t(0.05)$ and $C_t(0.50)$ that tracks how tail connectedness evolves through crisis episodes.


\bigskip\noindent\rule{\textwidth}{0.4pt}\bigskip


\paragraph{Module 4b: Tail-Risk Amplification (Quantile Regression with HLI Interaction)}

\textbf{Purpose}: Test whether high hidden leverage amplifies left-tail outcomes --- i.e., does the HLI predict \textit{worse} downside risk during stress?

\textbf{Unit of observation}: Quarter \$t\$ (aggregate time series).

\textbf{Input variables}:


{\small
\begin{longtable}{|p{0.15\textwidth}|p{0.15\textwidth}|p{0.15\textwidth}|p{0.15\textwidth}|p{0.15\textwidth}|p{0.15\textwidth}|}
\hline
\textbf{Variable} & \textbf{Symbol} & \textbf{Definition} & \textbf{Source} & \textbf{Frequency} & \textbf{Level} \\
\hline
\endhead
Outcome variable & $y_t$ & System-wide equity return or financial conditions index & Bloomberg / FRED & Quarterly & Aggregate \\
\hline
Stress indicator & $\text{Stress}_t$ & VIX level, or dummy for VIX \$>\$ 75th percentile & CBOE (via FRED) & Quarterly & Aggregate \\
\hline
Hidden Leverage Index & $\text{HLI}_t$ & From Module 3 & Module 3 output & Quarterly & Aggregate \\
\hline
Controls & $X_t$ & Term spread (10Y-2Y), credit spread (BAA-AAA), GDP growth & FRED & Quarterly & Aggregate \\
\hline
\end{longtable}
}

\textbf{Model specification} --- Quantile regression at $\tau = 0.05$:

\[Q_{\tau}(y_t) = \alpha(\tau) + \beta_1(\tau) \cdot \text{Stress}_t + \beta_2(\tau) \cdot \text{HLI}_{t-1} + \beta_3(\tau) \cdot (\text{Stress}_t \times \text{HLI}_{t-1}) + \gamma(\tau)' X_{t-1} + \varepsilon_t\]

Note: HLI is \textbf{lagged} (predictive, not contemporaneous) to address endogeneity.

\textbf{Key parameter}: $\beta_3(\tau = 0.05)$. If negative and significant: high hidden leverage makes left-tail outcomes \textit{worse} during stress. If $\beta_3(\tau = 0.50) \approx 0$: no effect at the median --- confirming this is a \textit{tail-specific} amplification mechanism.

\textbf{Estimation details}:
\begin{itemize}
\item \texttt{statsmodels.QuantReg} at $\tau \in \{0.05, 0.25, 0.50, 0.75, 0.95\}$.
\item Standard errors: Kernel (Hall-Sheather bandwidth) or bootstrap (1000 replications).
\item Formal interquantile test: Wald test of $H_0: \beta_3(0.05) = \beta_3(0.50)$.

\end{itemize}
\textbf{Output}: Table of coefficients across quantiles, showing that the HLI$\times$Stress interaction is large and negative at $\tau = 0.05$ but near zero at $\tau = 0.50$.


\bigskip\noindent\rule{\textwidth}{0.4pt}\bigskip


\paragraph{Module 4c: Growth-at-Risk (Adrian et al. 2019)}

\textbf{Purpose}: Test whether the HLI predicts worse left-tail GDP growth --- extending the GaR framework with an NBFI-specific predictor.

\textbf{Unit of observation}: Quarter \$t\$ (aggregate time series). Potentially also country \$c\$ $\times$ quarter \$t\$ (panel).

\textbf{Input variables}:


{\small
\begin{longtable}{|p{0.15\textwidth}|p{0.15\textwidth}|p{0.15\textwidth}|p{0.15\textwidth}|p{0.15\textwidth}|p{0.15\textwidth}|}
\hline
\textbf{Variable} & \textbf{Symbol} & \textbf{Definition} & \textbf{Source} & \textbf{Frequency} & \textbf{Level} \\
\hline
\endhead
Future GDP growth & $y_{t+h}$ & Real GDP growth at horizon \$h\$ (1, 4, 8, 12 quarters ahead) & FRED (GDPC1) & Quarterly & Aggregate \\
\hline
Financial conditions & $\text{FCI}_t$ & National Financial Conditions Index or composite (VIX + credit spread + term spread) & Chicago Fed NFCI; or constructed & Quarterly & Aggregate \\
\hline
Hidden Leverage Index & $\text{HLI}_t$ & From Module 3 & Module 3 output & Quarterly & Aggregate \\
\hline
\end{longtable}
}

\textbf{Model specification} --- Quantile regression:

\[Q_{\tau}(y_{t+h}) = \alpha(\tau) + \beta_1(\tau) \cdot \text{FCI}_t + \beta_2(\tau) \cdot \text{HLI}_t + \varepsilon_t\]

Estimated at multiple horizons $h \in \{1, 2, 4, 8, 12\}$ quarters and multiple quantiles $\tau \in \{0.05, 0.25, 0.50, 0.75, 0.95\}$.

\textbf{Key test}: $\beta_2(\tau = 0.05) < 0$ and significant, but $\beta_2(\tau = 0.50) \approx 0$. This means high hidden leverage predicts worse left-tail growth outcomes but has no effect on median growth --- confirming that shadow leverage is a \textit{tail risk} phenomenon.

\textbf{Estimation details}:
\begin{itemize}
\item \texttt{statsmodels.QuantReg} with \texttt{growth\_at\_risk()} function from \texttt{src/econometrics/quantile\_var.py}.
\item Also estimate the ''term structure of GaR'': how the FCI and HLI effects change across horizons \$h\$.
\item Extended sample: 2000Q1--2024Q4 using annual FSB data interpolated to quarterly frequency for the HLI.

\end{itemize}
\textbf{Output}:
\begin{enumerate}
\item \textbf{GaR term structure plot}: for each horizon \$h\$, the 5th percentile of predicted GDP growth, split by HLI high vs. HLI low.
\item \textbf{Coefficient table}: $\beta_2(\tau)$ across quantiles --- showing the HLI matters only in the left tail.
\item \textbf{Fan chart}: predicted distribution of future GDP growth conditional on current FCI and HLI values.

\end{enumerate}

\bigskip\noindent\rule{\textwidth}{0.4pt}\bigskip


\subsubsection{Module Linkage Map (Project 1)}


\begin{adjustbox}{max width=\textwidth}
\begin{minipage}{1.5\textwidth}
\begin{verbatim}
                    +----------------------------------+
  OFR / ECB repo    |  Module 1: Implied Leverage       |
  DTCC derivatives  |  Input: E, R, D per sub-sector    |    Output: Lev_{s,t}
  FSB monitoring    |  Formula: (E + R + deltaD) / E        |    panel
                    +--------------+-------------------+
                                   | Leverage weights edges
                                   v
                    +----------------------------------+
  BIS Locational    |  Module 2: Bipartite Network      |
  OFR bilateral     |  Input: Exp_{ij,t}, Lev_{j,t}    |    Output: W_t
  repo data         |  Construct: adjacency matrix W_t  |    adjacency matrices
                    +--------------+-------------------+
                                   | Spectral decomposition
                                   v
                    +----------------------------------+
                    |  Module 3: HLI (Spectral Index)   |
                    |  Input: W_t                       |    Output: HLI_t
                    |  Compute: lambda_1(W_t)                |    quarterly scalar
                    +--------------+-------------------+
                                   | Predictor in regressions
                      +------------+------------+
                      v            v            v
               +------------+ +----------+ +----------+
               |  4a: QVAR  | | 4b: Tail | | 4c: GaR  |
               |  Connect.  | | Amplif.  | |          |
               +------------+ +----------+ +----------+
\end{verbatim}
\end{minipage}
\end{adjustbox}


\subsection{Expected Contributions}

\begin{enumerate}
\item First implied leverage measure for NBFIs from publicly observable market data.
\item Hidden Leverage Index as a leading indicator --- spikes before March 2020 and September 2022.
\item Shadow leverage amplifies left-tail risk (GaR) but has no effect at the median.
\item Quantile connectedness confirms asymmetric contagion in bank-NBFI networks.

\end{enumerate}

\subsection{Data Sources \& Sample}


{\small
\begin{longtable}{|p{0.15\textwidth}|p{0.15\textwidth}|p{0.33\textwidth}|p{0.22\textwidth}|}
\hline
\textbf{Data Source} & \textbf{Variables} & \textbf{Frequency} & \textbf{Coverage} \\
\hline
\endhead
OFR US Repo (SOFR) & Repo volumes, rates, counterparty type & Daily & 2014--present \\
\hline
ECB MMSR & Euro repo volumes by counterparty & Daily & 2016--present \\
\hline
DTCC Swap Data Repository & IRS/CDS notional, counterparty type & Weekly & 2013--present \\
\hline
FSB Global Monitoring Report & NBFI sector AUM, leverage proxies & Annual & 2002--present \\
\hline
BIS Locational Banking Stats & Cross-border bank-NBFI exposures & Quarterly & 2000--present \\
\hline
FRED / ECB SDW & GDP, financial conditions, VIX, spreads & Mixed & 2000--present \\
\hline
\end{longtable}
}

\textbf{Sample period}: 2013Q1--2024Q4 (post-DTCC reporting). Extended to 2000 for GaR regressions using FSB annual data.

\textbf{Key stress episodes for validation}: Taper Tantrum (May 2013), China devaluation (Aug 2015), Gilt crisis (Sep 2022), COVID-19 (Mar 2020), SVB failure (Mar 2023).


\subsection{Literature Positioning}

This project contributes to three strands:

\begin{enumerate}
\item \textbf{NBFI leverage measurement}: Extends Jiang, Matvos, Piskorski \& Seru (2024) on hidden bank losses to the NBFI sector. Complements FSB (2023) Global Monitoring Report with a market-data-based approach.

\item \textbf{Spectral network measures}: Builds on Acemoglu, Ozdaglar \& Tahbaz-Salehi (2015) and Greenwood, Landier \& Thesmar (2015) on network-based contagion measures. The HLI adapts the spectral radius concept to a bipartite bank-NBFI setting.

\item \textbf{Growth-at-Risk}: Extends Adrian, Boyarchenko \& Giannone (2019) by adding NBFI-specific predictors (hidden leverage) to the financial conditions $\rightarrow$ GDP growth-at-risk framework.

\end{enumerate}

\subsection{Identification \& Robustness}

\begin{itemize}
\item \textbf{Endogeneity of leverage}: Leverage is both a response to and a predictor of stress. We address this via (a) lagged HLI values, (b) Granger causality tests, and (c) instrumenting leverage with regulatory threshold dummies (e.g., margin call triggers).
\item \textbf{Measurement error in implied leverage}: Delta-equivalence factor $\delta$ is calibrated from BIS survey data; robustness across $\delta \in [0.02, 0.10]$.
\item \textbf{Network construction sensitivity}: Results tested across alternative edge-weighting schemes (gross vs net, bilateral vs multilateral netting).

\end{itemize}

\bigskip\noindent\rule{\textwidth}{0.4pt}\bigskip


\section{Project 2: Mapping the Channels}

\textbf{Subtitle:} Network Analysis of Non-Bank Amplification in Monetary Policy Transmission


\subsection{Motivation}

Building on Banerjee, Hofmann, Ng \& Pinter (2025, BIS Bulletin 116), who document
that other financial intermediaries (OFIs) amplify monetary policy transmission to
long-term yields and credit spreads while insurance companies and pension funds
(ICPFs) dampen it, we decompose this aggregate effect using granular network and
connectedness methods. Their panel local-projection analysis reveals significant
amplification but with wide confidence bands, leaving open: \textit{which} OFI sub-sectors
drive it, \textit{through which} balance-sheet channels, and \textit{under what} conditions.


\subsection{Hypotheses}

\textbf{H1 (Heterogeneous amplification).} Procyclicality beta is highest for hedge
funds and leveraged investment funds, moderate for bond mutual funds and MMFs,
and lowest for pension funds and insurance.

\textbf{H2 (Nonlinear tail transmission).} Quantile IRFs at $\tau = 0.05$ show 2--3$\times$
stronger responses than at $\tau = 0.50$, especially for the policy rate $\rightarrow$ hedge
fund returns link.

\textbf{H3 (Variance amplification).} NBFI leverage increases the conditional variance
(not just the conditional mean) of future growth outcomes, generating fatter
left tails.

\textbf{H4 (Network containment).} During the 2022--23 hiking cycle, within-OFI
connectedness was high but OFI-to-real-economy connectedness remained low,
explaining the resilience puzzle.


\subsection{Empirical Strategy}

Project 2 has \textbf{five modules}. Module 1 establishes the micro-foundation (VaR-constraint channel). Modules 2--3 provide the econometric core. Modules 4--5 address variance effects and the 2022--23 puzzle.


\begin{adjustbox}{max width=\textwidth}
\begin{minipage}{1.5\textwidth}
\begin{verbatim}
Module 1 (VaR Channel)  --→  Module 2 (DCC-GARCH)  --→  Module 3 (Quantile VAR)  --→  Module 4 (LS-GaR)  --→  Module 5 (Banerjee Puzzle)
       ↓                            ↓                           ↓                          ↓                        ↓
  Sub-sector beta_s            Time-varying bank-         Tail-specific MP              NBFI leverage         Within-OFI vs
  procyclicality            NBFI correlations          transmission IRFs             fattens tails         OFI-to-real
  + fire-sale sim           (regime detection)         + connectedness               (variance channel)    economy split
\end{verbatim}
\end{minipage}
\end{adjustbox}


\bigskip\noindent\rule{\textwidth}{0.4pt}\bigskip


\subsubsection{Module 1: VaR-Constraint Channel (Adrian \& Shin 2010, 2014)}


\paragraph{Economic logic}

Financial intermediaries subject to Value-at-Risk (VaR) constraints manage their balance sheets by targeting a maximum loss at a given confidence level. When asset prices rise, measured portfolio risk (VaR) falls mechanically, creating unused risk capacity. Intermediaries respond by expanding their balance sheets --- buying more assets funded by short-term debt --- so that leverage \textit{rises} with total assets. This is the \textbf{procyclicality of leverage}: $\text{Corr}(\Delta \log \text{Leverage}, \Delta \log \text{Assets}) > 0$.

The reverse is more dangerous. When asset prices fall, VaR rises, breaching the constraint. Intermediaries must sell assets to restore compliance, which depresses prices further, triggering more VaR breaches across interconnected institutions --- the \textbf{fire-sale spiral}. Adrian \& Shin (2010, 2014) document this mechanism for US broker-dealers. We extend it to the full NBFI universe and test whether the mechanism is heterogeneous across sub-sectors.


\paragraph{What we test}

\begin{enumerate}
\item \textbf{Procyclicality by sub-sector (H1)}: Is leverage procyclical for each NBFI type, and is the procyclicality coefficient $\beta_s$ larger for leveraged sub-sectors (hedge funds) than for liability-driven ones (pension funds, insurance)?

\item \textbf{Fire-sale amplification}: Does adding NBFI sub-sectors to a bank-only fire-sale simulation produce larger aggregate losses, and which sub-sectors contribute most to the amplification?

\end{enumerate}

\bigskip\noindent\rule{\textwidth}{0.4pt}\bigskip


\paragraph{Specification 1: Procyclicality of leverage (sub-sector panel)}

\textbf{Unit of observation}: NBFI sub-sector \$s\$ in quarter \$t\$.

\textbf{Input variables}:


{\small
\begin{longtable}{|p{0.15\textwidth}|p{0.15\textwidth}|p{0.15\textwidth}|p{0.15\textwidth}|p{0.15\textwidth}|p{0.15\textwidth}|}
\hline
\textbf{Variable} & \textbf{Symbol} & \textbf{Definition} & \textbf{Source} & \textbf{Frequency} & \textbf{Level} \\
\hline
\endhead
Leverage & $\text{Lev}_{s,t}$ & Total assets / equity (net asset value) of sub-sector \$s\$ & Fed Z.1 Flow of Funds; ECB Investment Fund Statistics & Quarterly & Sub-sector aggregate \\
\hline
Total assets & $A_{s,t}$ & Total financial assets of sub-sector \$s\$ & Fed Z.1 (tables L.122--L.130) & Quarterly & Sub-sector aggregate \\
\hline
Equity & $E_{s,t}$ & Net asset value (= assets - liabilities) of sub-sector \$s\$ & Fed Z.1 & Quarterly & Sub-sector aggregate \\
\hline
VIX & $\text{VIX}_t$ & CBOE Volatility Index (quarterly average of daily values) & FRED (VIXCLS) & Daily $\rightarrow$ Quarterly avg & Aggregate \\
\hline
Term spread & --- & 10Y Treasury - 2Y Treasury yield & FRED (T10Y2Y) & Daily $\rightarrow$ Quarterly avg & Aggregate \\
\hline
Credit spread & --- & BAA - AAA corporate bond spread & FRED (BAA - AAA) & Daily $\rightarrow$ Quarterly avg & Aggregate \\
\hline
GDP growth & --- & Real GDP quarter-on-quarter growth (annualised) & FRED (GDPC1) & Quarterly & Aggregate \\
\hline
\end{longtable}
}

\textbf{Sub-sectors} $s \in \{$hedge funds, MMFs, bond mutual funds, equity mutual funds, ETFs, insurance companies, pension funds, finance companies, securitisation vehicles$\}$.

\textbf{Model specification}:

\[\Delta \log(\text{Lev}_{s,t}) = \alpha_s + \beta_s \cdot \Delta \log(A_{s,t}) + \gamma \cdot X_{s,t-1} + \delta_t + \varepsilon_{s,t}\]

where:
\begin{itemize}
\item $\text{Lev}_{s,t} = A_{s,t} / E_{s,t}$ is the leverage ratio of sub-sector \$s\$
\item $X_{s,t-1}$ = lagged controls: VIX level, term spread, credit spread, GDP growth
\item $\alpha_s$ = \textbf{sub-sector fixed effects} (absorb time-invariant differences in business models, regulatory treatment, and average leverage levels)
\item $\delta_t$ = \textbf{time fixed effects} (absorb common macroeconomic shocks --- monetary policy, global risk appetite --- that affect all sub-sectors simultaneously)
\item $\beta_s$ = sub-sector-specific procyclicality coefficient (the \textbf{key parameter of interest})

\end{itemize}
\textbf{Interpretation of $\beta_s$}: $\beta_s > 0$ means leverage is procyclical (rises when assets grow, falls when assets shrink) --- consistent with active VaR-constraint management. $\beta_s \approx 0$ means leverage is roughly constant (passive balance sheet). $\beta_s < 0$ would mean counter-cyclical leverage.

\textbf{Expected ranking} (H1): $\beta_{\text{HF}} > \beta_{\text{Lev. funds}} > \beta_{\text{Bond MF}} > \beta_{\text{MMF}} > \beta_{\text{IC}} \approx \beta_{\text{PF}} \approx 0$.

\textbf{Estimation details}:
\begin{itemize}
\item OLS panel regression with entity and time fixed effects.
\item Standard errors clustered at the sub-sector level (9 clusters). Robustness: two-way clustering (sub-sector $\times$ year).
\item In code: \texttt{sector\_procyclicality()} in \texttt{src/models/var\_amplification.py} estimates this for each sub-sector using rolling volatility as a VaR-implied leverage proxy: $\text{Lev}_{\text{implied}} = 1/(z_\alpha \cdot \sigma)$ where $z_\alpha$ is the VaR critical value and $\sigma$ is rolling volatility.

\end{itemize}
\textbf{Output}: A table of $\beta_s$ coefficients with standard errors, one row per sub-sector. Bar chart showing the ranking.


\bigskip\noindent\rule{\textwidth}{0.4pt}\bigskip


\paragraph{Specification 2: Interaction with stress}

\textbf{Purpose}: Test whether procyclicality \textit{intensifies} during stress (asymmetric amplification).

\textbf{Model specification}:

\[\Delta \log(\text{Lev}_{s,t}) = \alpha_s + \beta_1 \cdot \Delta \log(A_{s,t}) + \beta_2 \cdot \Delta \log(A_{s,t}) \times \text{Stress}_t + \gamma \cdot X_{s,t-1} + \delta_t + \varepsilon_{s,t}\]

where $\text{Stress}_t$ = dummy for quarters with VIX \$>\$ 75th percentile (or NBER recession indicator).

\textbf{Key test}: If $\beta_2 > 0$: procyclicality is \textit{stronger} during stress --- exactly when it does most damage.


\bigskip\noindent\rule{\textwidth}{0.4pt}\bigskip


\paragraph{Specification 3: Fire-sale simulation (Greenwood, Landier \& Thesmar 2015)}

\textbf{Purpose}: Quantify how much NBFIs amplify fire-sale losses beyond a bank-only system.

\textbf{Unit of observation}: Simulated system (agent-based, not regression).

\textbf{Input variables (simulation parameters)}:


{\small
\begin{longtable}{|p{0.15\textwidth}|p{0.15\textwidth}|p{0.33\textwidth}|p{0.22\textwidth}|}
\hline
\textbf{Parameter} & \textbf{Symbol} & \textbf{Definition} & \textbf{Value} \\
\hline
\endhead
Initial shock & --- & Exogenous decline in asset prices & -5\% \\
\hline
Market depth & $\lambda^{-1}$ & Kyle (1985) price impact parameter (\$ volume needed to move price 1\%) & \\$500K (calibrated) \\
\hline
VaR confidence & $z_\alpha$ & Confidence level for VaR constraint & 99\% (banks), 95--99.5\% (NBFIs) \\
\hline
GARCH decay & --- & EWMA lambda for volatility updating & 0.94 \\
\hline
Bank system & --- & 10 banks, leverage 10--15$\times$, $\sigma$ 1--2\% & Calibrated to G-SIB data \\
\hline
NBFI system & --- & 20 entities: 5 HFs (leverage 5--25$\times$), 8 investment funds (1--3$\times$), 4 pension funds (1--2$\times$), 3 insurance (3--8$\times$) & Calibrated to FSB data \\
\hline
\end{longtable}
}

\textbf{Simulation mechanism} (multi-round):
\begin{enumerate}
\item \textbf{Initial shock}: Asset prices fall by -5\%.
\item \textbf{Volatility update}: Each entity updates its rolling volatility estimate (EWMA, $\lambda = 0.94$).
\item \textbf{VaR tightening}: Higher volatility means the VaR constraint binds more --- entities with VaR-constrained leverage ($\beta_s > 0$) must reduce exposure.
\item \textbf{Forced selling}: Entities sell assets to restore target leverage. Selling volume = excess leverage $\times$ assets.
\item \textbf{Price impact}: Aggregate selling depresses prices via Kyle lambda: $\Delta p = -\text{TotalSelling} / \text{MarketDepth}$.
\item \textbf{Repeat}: Steps 2--5 iterate until convergence (total selling $< 0.01\%$ of assets) or max 50 rounds.

\end{enumerate}
\textbf{Comparison}:
\begin{itemize}
\item \textit{Bank-only system}: Only 10 banks participate in fire-sale rounds.
\item \textit{Bank + NBFI system}: All 30 entities participate with their calibrated parameters.
\item \textbf{Amplification ratio} = Total price drop (bank+NBFI) / Total price drop (bank-only). If \$> 1\$: NBFIs amplify.

\end{itemize}
\textbf{Estimation details}:
\begin{itemize}
\item Implemented as \texttt{FireSaleSimulation} dataclass in \texttt{src/models/var\_amplification.py}.
\item \texttt{run\_counterfactual()} runs both systems with identical initial shock and compares.
\item \texttt{amplification\_ratio()} computes the ratio plus equity losses for each system.

\end{itemize}
\textbf{Output}:
\begin{enumerate}
\item \textbf{Amplification ratio}: e.g., 1.8$\times$ means bank+NBFI system suffers 80\% larger price drop.
\item \textbf{Round-by-round dynamics}: time series of asset prices, aggregate leverage, total equity, and number of active entities --- for both systems.
\item \textbf{3-panel comparison chart}: asset price path, leverage path, and equity path for bank-only vs. bank+NBFI.

\end{enumerate}
\textbf{$\rightarrow$ Link to Module 2}: The procyclicality coefficients $\beta_s$ from Specification 1 inform which sub-sectors are most dangerous in the fire-sale simulation. Module 2 (DCC-GARCH) tests whether the bank-NBFI correlations that drive fire-sale contagion are indeed time-varying and stress-dependent.


\bigskip\noindent\rule{\textwidth}{0.4pt}\bigskip


\subsubsection{Module 2: DCC-GARCH Dynamic Correlations (Engle 2002)}

\textbf{Purpose}: Estimate time-varying pairwise correlations between bank and NBFI sector returns, testing whether correlations spike during stress (confirming that contagion is regime-dependent, not constant).

\textbf{Unit of observation}: Pair of return series (bank \$i\$, NBFI sub-sector \$j\$) $\times$ week \$t\$.

\textbf{Input variables}:


{\small
\begin{longtable}{|p{0.15\textwidth}|p{0.15\textwidth}|p{0.15\textwidth}|p{0.15\textwidth}|p{0.15\textwidth}|p{0.15\textwidth}|}
\hline
\textbf{Variable} & \textbf{Symbol} & \textbf{Definition} & \textbf{Source} & \textbf{Frequency} & \textbf{Level} \\
\hline
\endhead
Bank returns & $r_{i,t}$ & Weekly log-returns for individual G-SIBs or bank sector index & Bloomberg / Yahoo Finance & Weekly & Institution \\
\hline
NBFI returns & $r_{j,t}$ & Weekly log-returns for NBFI sub-sector indices or representative firms & Bloomberg (BLK, BX, KKR, MET, PRU, etc.) & Weekly & Institution or sub-sector \\
\hline
\end{longtable}
}

\textbf{Typical system}: 6 entities --- e.g., top 3 banks (JPM, GS, HSBC) + top 3 NBFIs (BLK, BX, MET). Or sector-level indices.

\textbf{Model specification} --- Two-step DCC(1,1):

\textbf{Step 1 --- Univariate GARCH(1,1)} for each series \$i\$:

\[r_{i,t} = \mu_i + \varepsilon_{i,t}, \quad \varepsilon_{i,t} = \sigma_{i,t} z_{i,t}, \quad z_{i,t} \sim N(0,1)\]

\[\sigma_{i,t}^2 = \omega_i + \alpha_i \varepsilon_{i,t-1}^2 + \beta_i \sigma_{i,t-1}^2\]

This gives: conditional volatilities $\sigma_{i,t}$ and standardised residuals $z_{i,t} = \varepsilon_{i,t}/\sigma_{i,t}$.

\textbf{Step 2 --- DCC dynamics} on the standardised residuals:

\[Q_t = (1 - a - b)\bar{Q} + a(z_{t-1}z_{t-1}') + b Q_{t-1}\]

\[R_t = \text{diag}(Q_t)^{-1/2} \cdot Q_t \cdot \text{diag}(Q_t)^{-1/2}\]

where:
\begin{itemize}
\item $\bar{Q}$ = unconditional correlation matrix of $z_t$ (estimated from full sample)
\item \$a\$ = ''news'' parameter --- how fast correlations react to shocks (typically 0.01--0.10)
\item \$b\$ = ''persistence'' parameter --- how slowly correlations revert (typically 0.85--0.98)
\item $R_t$ = time-varying correlation matrix (the object of interest)

\end{itemize}
\textbf{Estimation details}:
\begin{itemize}
\item \textbf{Step 1}: \texttt{arch} package in Python. \texttt{fit\_univariate\_garch()} in \texttt{src/models/dcc\_garch.py}. GARCH(1,1) with constant mean. Standard MLE.
\item \textbf{Step 2}: Quasi-MLE on DCC log-likelihood. \texttt{estimate\_dcc()} uses Nelder-Mead optimisation with initial values $a_0 = 0.05$, $b_0 = 0.90$. Constraint: \$a > 0\$, \$b > 0\$, \$a + b < 1\$.
\item Persistence \$= a + b\$. High persistence (e.g., \$0.98\$) means correlations are very slowly mean-reverting.

\end{itemize}
\textbf{Output}:
\begin{enumerate}
\item \textbf{Dynamic pairwise correlations} $\rho_{ij,t}$: time series of correlation between every pair of entities. Plotted as line charts over 2000--2024.
\item \textbf{Sector-average correlations}: \texttt{sector\_average\_correlation()} averages all bank-HF pairs, all bank-insurance pairs, etc. Shows which bank-NBFI links are tightest.
\item \textbf{Contagion test} (Forbes \& Rigobon 2002): \texttt{correlation\_breakdown\_test()} formally tests whether correlations are \textit{significantly higher} during crisis periods (VIX > 75th percentile) vs. calm periods. Output: mean correlation in each regime, difference, t-statistic, p-value.

\end{enumerate}
\textbf{Key findings expected}:
\begin{itemize}
\item Bank-hedge fund correlations spike during stress (confirming fire-sale contagion channel from Module 1).
\item Bank-pension fund correlations are low and stable (confirming that liability-driven entities don't amplify).
\item Bank-insurance correlations increase moderately (mixed evidence).

\end{itemize}
\textbf{$\rightarrow$ Link to Module 3}: DCC-GARCH captures \textit{mean} correlation dynamics. Module 3 (Quantile VAR) goes further by estimating \textit{tail-specific} dependence --- testing whether the system is more connected at the 5th percentile than at the median.


\bigskip\noindent\rule{\textwidth}{0.4pt}\bigskip


\subsubsection{Module 3: Quantile VAR --- Nonlinear Transmission (Core Innovation)}

\textbf{Purpose}: Estimate how monetary policy shocks propagate through the bank-NBFI system at different points of the return distribution. The key question: is MP transmission a \textit{tail phenomenon} --- much stronger at the 5th percentile than at the median?

\textbf{Unit of observation}: System of \$k\$ return/macro series $\times$ quarter \$t\$ (or month \$t\$).

\textbf{Input variables}:


{\small
\begin{longtable}{|p{0.15\textwidth}|p{0.15\textwidth}|p{0.15\textwidth}|p{0.15\textwidth}|p{0.15\textwidth}|p{0.15\textwidth}|}
\hline
\textbf{Variable} & \textbf{Symbol} & \textbf{Definition} & \textbf{Source} & \textbf{Frequency} & \textbf{Level} \\
\hline
\endhead
Policy rate & $r_t^{MP}$ & Effective Fed funds rate (or shadow rate during ZLB) & FRED (DFF; Wu-Xia shadow rate) & Monthly $\rightarrow$ Quarterly & Aggregate \\
\hline
Bank sector return & $r_t^{B}$ & Aggregate bank equity return index & Bloomberg (KBW Bank Index or constructed) & Weekly $\rightarrow$ Quarterly & Sector \\
\hline
Hedge fund return & $r_t^{HF}$ & Hedge fund composite return & HFRX Global Index; Bloomberg & Monthly $\rightarrow$ Quarterly & Sector \\
\hline
MMF return & $r_t^{MMF}$ & Money market fund return (proxy: short-term yield changes) & iMoneyNet; FRED (money market yields) & Monthly $\rightarrow$ Quarterly & Sector \\
\hline
Bond fund return & $r_t^{BF}$ & Bond mutual fund return & Bloomberg (AGG total return) & Monthly $\rightarrow$ Quarterly & Sector \\
\hline
Insurance return & $r_t^{IC}$ & Insurance sector equity return & Bloomberg (S\&P Insurance Index) & Weekly $\rightarrow$ Quarterly & Sector \\
\hline
Credit spread & $s_t$ & BAA--AAA corporate bond spread & FRED (BAA - AAA) & Monthly $\rightarrow$ Quarterly & Aggregate \\
\hline
Term spread & --- & 10Y - 2Y Treasury & FRED (T10Y2Y) & Monthly $\rightarrow$ Quarterly & Aggregate \\
\hline
\end{longtable}
}

\textbf{System vector}: $Y_t = (r_t^{MP}, r_t^{B}, r_t^{HF}, r_t^{MMF}, r_t^{BF}, r_t^{IC}, s_t)'$ --- a \textbf{7-variable Quantile VAR}.

The choice of variables is driven by the economic question: we want to trace how a monetary policy shock ($r^{MP}$) propagates to each NBFI sub-sector return and to credit spreads, and whether this transmission is stronger in the tails.

\textbf{Model specification} --- Quantile VAR(p) at quantile $\tau$:

\[Q_{\tau}(y_{i,t} \mid Y_{t-1}, \ldots, Y_{t-p}) = c_i(\tau) + \sum_{j=1}^{7} \sum_{l=1}^{p} \beta_{ij}^{(l)}(\tau) \, y_{j,t-l}\]

For each of the 7 equations, the dependent variable $y_{i,t}$ is regressed on \$p\$ lags of all 7 variables. This gives $7 \times (7p + 1)$ coefficients \textit{at each quantile}.


\paragraph{What the QVAR captures (and what it does not)}

A standard VAR models the conditional \textbf{mean} of each variable: $E[y_{i,t} | Y_{t-1}]$. The coefficients are the same regardless of whether we are in calm or stressed times. This is a strong assumption --- it says the transmission mechanism is symmetric.

The Quantile VAR replaces the conditional mean with the conditional \textbf{quantile}: $Q_\tau(y_{i,t} | Y_{t-1})$. Now \textbf{every coefficient depends on $\tau$}. Critically, this does \textit{not} mean ''where is the regressor in its distribution.'' It means: \textbf{what predicts the $\tau$-th quantile of the dependent variable?}

\begin{itemize}
\item At $\tau = 0.50$: you model the \textbf{median} of bank returns --- ''normal times'' transmission.
\item At $\tau = 0.05$: you model the \textbf{5th percentile} (left tail) of bank returns --- ''how do other variables affect the worst outcomes for banks?''

\end{itemize}
So $\beta_{ij}^{(l)}(\tau = 0.05)$ tells you how a one-unit change in $y_j$ at lag \$l\$ shifts the 5th percentile of $y_i$. This can be completely different from $\beta_{ij}^{(l)}(\tau = 0.50)$. For instance, a monetary policy tightening might have a small effect on the median of hedge fund returns but a devastating effect on the 5th percentile --- because only the worst-performing funds face margin calls and forced liquidation.


\paragraph{The Ando, Greenwood-Nimmo \& Shin (2022) framework}

Once the QVAR is estimated at quantile $\tau$, the companion matrix gives the MA($\infty$) representation, from which we compute a Generalised Forecast Error Variance Decomposition (GFEVD) --- \textbf{at each quantile separately}. This is the key contribution of Ando et al. (2022): applying the Diebold-Yilmaz (2012) connectedness framework not at the mean but across the quantile distribution.

The GFEVD entry $\theta_{ij}(\tau)$ answers: ''What share of variable \$i\$'s \$h\$-step-ahead forecast error variance (at quantile $\tau$) is attributable to shocks originating in variable \$j\$?'' From this $k \times k$ matrix $\Theta(\tau)$, Ando et al. define:
\begin{itemize}
\item \textbf{Total connectedness} $C(\tau)$: the average off-diagonal share --- how much of each variable's risk comes from \textit{other} variables.
\item \textbf{Directional measures}: TO (how much \$i\$ transmits), FROM (how much \$i\$ receives), NET (transmitter vs. receiver).

\end{itemize}
The critical insight: \textbf{you compute $C(\tau)$ at every quantile and compare}. If $C(0.05) \gg C(0.50)$, it means the system is dramatically more interconnected when variables are in their left tails. In normal times, shocks to one sector explain perhaps 30\% of variation in others. In stress, they might explain 70\%. This asymmetry is invisible to a standard linear VAR --- which gives you one number regardless of whether markets are calm or crashing.


\paragraph{State-dependence and distance to constraints}

An important refinement: shock propagation depends not only on structural integration (who is connected to whom) but also on \textbf{how close intermediaries are to their VaR/margin/redemption constraints} when the shock hits. A moderate monetary policy tightening in a world where hedge funds have plenty of excess margin capacity is very different from the same shock when funds are already near their VaR limits.

The QVAR captures this implicitly: the coefficients $\beta_{ij}(\tau)$ vary with $\tau$ precisely because the transmission mechanism changes when the system is stressed. But we can make the state-dependence explicit with a \textbf{threshold specification}:

\[Q_\tau(y_{i,t}) = \begin{cases} \alpha_1(\tau) + \beta_1(\tau) X_t & \text{if } \text{Slack}_t > \bar{s} \text{ (far from constraints)} \\ \alpha_2(\tau) + \beta_2(\tau) X_t & \text{if } \text{Slack}_t \leq \bar{s} \text{ (near constraints)} \end{cases}\]

where $\text{Slack}_t$ measures the distance to binding constraints (e.g., aggregate excess margin, or the gap between actual VaR and VaR limit). The hypothesis: $\beta_2(\tau = 0.05) \gg \beta_1(\tau = 0.05)$ --- even in the left tail, amplification is much worse when constraints are nearly binding. This is a \textbf{double nonlinearity}: tail-specific \textit{and} state-dependent.

This connects to the HLI from Project 1: high $\lambda_1(W_t)$ (tight network) \textbf{plus} low slack (near constraints) is the dangerous combination. Neither alone is sufficient.

\textbf{Estimation details}:
\begin{itemize}
\item \textbf{Equation-by-equation quantile regression}: Each equation estimated separately via \texttt{statsmodels.QuantReg}. This is computationally feasible even for 7 equations because each is a standard linear quantile regression.
\item \textbf{Lags}: \$p = 4\$ for quarterly data (one year of history). BIC-optimal at each quantile; robustness with $p \in \{1, 2, 4\}$.
\item \textbf{Quantiles}: $\tau \in \{0.05, 0.25, 0.50, 0.75, 0.95\}$.
\item \textbf{Residuals}: After estimation, compute residuals $\hat{u}_{i,t}(\tau) = y_{i,t} - \hat{Q}_\tau(y_{i,t}|\cdot)$ and their covariance matrix $\hat{\Sigma}(\tau)$.
\item Implemented in \texttt{estimate\_quantile\_var()} and \texttt{estimate\_quantile\_var\_grid()} in \texttt{src/econometrics/quantile\_var.py}.

\end{itemize}
\textbf{From the QVAR, three outputs are computed:}

\textbf{(a) Quantile Impulse Response Functions (QIRFs)}:
\begin{itemize}
\item Trace the effect of a 1-standard-deviation shock to variable \$j\$ (e.g., policy rate) on variable \$i\$ (e.g., hedge fund return) over \$h\$ periods.
\item Uses the companion matrix representation: shock is propagated through the MA($\infty$) representation.
\item Computed via \texttt{quantile\_irf()}. Key comparison: \texttt{compare\_quantile\_irfs()} shows the IRF at $\tau = 0.05$ vs. $\tau = 0.50$.
\item \textbf{Expected result}: A monetary policy tightening shock has a 2--3$\times$ larger negative effect on hedge fund returns at $\tau = 0.05$ than at $\tau = 0.50$.

\end{itemize}
\textbf{(b) Quantile Connectedness} (Ando, Greenwood-Nimmo \& Shin 2022):
\begin{itemize}
\item From the QVAR, compute the GFEVD $\Theta(\tau)$ at horizon \$h = 10\$.
\item Total connectedness: $C(\tau) = \frac{1}{k}\sum_{i \neq j} \theta_{ij}(\tau) \times 100\%$.
\item Directional: TO, FROM, NET for each variable.
\item Computed via \texttt{quantile\_connectedness()} in \texttt{src/econometrics/quantile\_var.py}.
\item \textbf{Expected result}: $C(0.05) \gg C(0.50)$ --- the system is much more interconnected in the left tail.

\end{itemize}
\textbf{(c) Rolling Quantile Connectedness}:
\begin{itemize}
\item \texttt{rolling\_quantile\_connectedness(window=60, step=1)} produces a time series $C_t(\tau)$ at each quantile.
\item Tracks how tail connectedness evolves over 2000--2024, spiking around crisis episodes.

\end{itemize}
\textbf{Output summary}:


{\small
\begin{longtable}{|p{0.22\textwidth}|p{0.33\textwidth}|p{0.33\textwidth}|}
\hline
\textbf{Output} & \textbf{What it shows} & \textbf{Format} \\
\hline
\endhead
QIRF comparison plot & MP shock $\rightarrow$ HF return at $\tau = 0.05$ vs \$0.50\$ & Line chart (horizon on x-axis) \\
\hline
Connectedness table & $C(\tau)$ at each quantile & Single row: \$C(0.05), C(0.25), C(0.50), C(0.75), C(0.95)\$ \\
\hline
GFEVD heatmap & $\Theta(\tau)$ matrix at $\tau = 0.05$ vs \$0.50\$ & $7 \times 7$ heatmap (two panels) \\
\hline
Directional bar chart & TO, FROM, NET by variable at each quantile & Grouped bar chart \\
\hline
Rolling connectedness & $C_t(0.05)$ and $C_t(0.50)$ over time & Time series (2 lines) \\
\hline
\end{longtable}
}


\paragraph{Extracting the Global Tail Cycle (GTC)}

The rolling quantile connectedness $C_t(0.05)$ captures how tightly the system is coupled \textit{in the tails} at each point in time. But Project 3 also needs a \textbf{common tail factor} --- a single series that captures when global markets are jointly in their left tails. This is the Global Tail Cycle (GTC), and it is a direct by-product of the QVAR estimation.

\textbf{Construction} (two approaches, both from the QVAR output):

\begin{enumerate}
\item \textbf{Quantile factor model (Ando \& Bai 2020)}: Estimate a quantile factor model at $\tau = 0.05$ on the \$k\$-variable system. The first common factor $\hat{f}_t(0.05)$ is the GTC --- the latent variable that best explains co-movement in the 5th percentile across all sectors. This is conceptually analogous to the Global Financial Cycle (Rey 2013), which is the first principal component at the mean, but extracted at the left tail.
\item \textbf{QVAR residual approach}: From the estimated QVAR at $\tau = 0.05$, compute the generalised forecast error variance decomposition. The first principal component of the off-diagonal spillover matrix $\Theta_t(0.05)$ gives a ''tail spillover factor'' that tracks when cross-sector tail transmission is elevated.

\end{enumerate}
\textbf{Output}: A quarterly time series $\text{GTC}_t$ exported for use in Project 3 (Module 3: NBFI amplification regressions) and Project 4 (as an alternative to the GFC for shock identification). The GTC differs from the GFC in that it captures co-movement specifically in crash episodes --- markets can appear diversified on average ($C_t(0.50)$ low) while being tightly coupled in the tails ($C_t(0.05)$ high).

\textbf{$\rightarrow$ Link to Module 4}: Module 3 shows that transmission is stronger in the tails. Module 4 tests whether NBFI leverage explains \textit{why} --- via the variance (scale) channel, not just the mean (location) channel.


\bigskip\noindent\rule{\textwidth}{0.4pt}\bigskip


\subsubsection{Module 4: Location-Scale Growth-at-Risk}

\textbf{Purpose}: Test whether NBFI leverage increases the \textit{conditional variance} (not just the conditional mean) of future growth outcomes, generating fatter left tails. This is the ''variance amplification'' hypothesis (H3).

\textbf{Unit of observation}: Quarter \$t\$ (aggregate time series) or country \$c\$ $\times$ quarter \$t\$ (panel).

\textbf{Input variables}:


{\small
\begin{longtable}{|p{0.15\textwidth}|p{0.15\textwidth}|p{0.15\textwidth}|p{0.15\textwidth}|p{0.15\textwidth}|p{0.15\textwidth}|}
\hline
\textbf{Variable} & \textbf{Symbol} & \textbf{Definition} & \textbf{Source} & \textbf{Frequency} & \textbf{Level} \\
\hline
\endhead
Future GDP growth & $y_{t+h}$ & Real GDP growth at horizon \$h\$ quarters ahead & FRED (GDPC1) & Quarterly & Aggregate \\
\hline
Financial conditions & $\text{FCI}_t$ & Chicago Fed NFCI or composite index & FRED (NFCI) & Quarterly & Aggregate \\
\hline
NBFI leverage & $\text{NBFI}_t$ & Aggregate NBFI leverage (from P1 Module 1), or NBFI assets as \% of GDP (FSB) & Module 1 output; FSB & Quarterly / Annual & Aggregate \\
\hline
Interaction & $\text{FCI}_t \times \text{NBFI}_t$ & Tests whether the FCI effect on variance depends on NBFI size & Constructed & Quarterly & Aggregate \\
\hline
\end{longtable}
}

\textbf{Model specification} --- Two-step Location-Scale model:

\textbf{Step 1 --- Location (conditional mean):}

\[\hat{\mu}_{t+h} = E[y_{t+h} \mid X_t] = \alpha + \beta_1 \cdot \text{FCI}_t + \beta_2 \cdot \text{NBFI}_t + \beta_3 \cdot (\text{FCI}_t \times \text{NBFI}_t)\]

Estimated by OLS. Residuals: $\hat{e}_{t+h} = y_{t+h} - \hat{\mu}_{t+h}$.

\textbf{Step 2 --- Scale (conditional variance):}

\[\log(\hat{e}_{t+h}^2) = \delta_0 + \delta_1 \cdot \text{FCI}_t + \delta_2 \cdot \text{NBFI}_t + \delta_3 \cdot (\text{FCI}_t \times \text{NBFI}_t) + v_t\]

Estimated by OLS on the log-squared residuals from Step 1. This gives $\hat{\sigma}_{t+h}^2 = \exp(\hat{\delta}_0 + \hat{\delta}_1 \text{FCI}_t + \ldots)$.

\textbf{Conditional quantiles} are then:

\[\hat{Q}_\tau(y_{t+h} \mid X_t) = \hat{\mu}_{t+h} + \hat{\sigma}_{t+h} \cdot \Phi^{-1}(\tau)\]

where $\Phi^{-1}(\tau)$ is the standard normal quantile function (or Student-t quantile for heavier tails).

\textbf{Key parameters}:
\begin{itemize}
\item $\delta_2 > 0$: NBFI leverage \textit{increases conditional variance} of future growth (fatter tails in both directions).
\item $\delta_3 > 0$: NBFI leverage amplifies the variance-increasing effect of tight financial conditions.
\item Standard GaR (Adrian et al. 2019) only models the location; our LS extension separates \textit{mean shift} from \textit{variance amplification}.

\end{itemize}
\textbf{Estimation details}:
\begin{itemize}
\item Implemented in \texttt{location\_scale\_gar()} in \texttt{src/econometrics/location\_scale.py}.
\item Step 1: OLS with HC1 standard errors.
\item Step 2: OLS on log-squared residuals.
\item Alternative: \texttt{skewt\_location\_scale\_fit()} estimates both location and scale parameters jointly via MLE assuming a Hansen (1994) skewed-t distribution (captures asymmetry).
\item Horizons: $h \in \{1, 4, 8, 12\}$ quarters.
\item Quantiles computed: $\tau \in \{0.05, 0.10, 0.25, 0.50, 0.75, 0.90, 0.95\}$.

\end{itemize}
\textbf{Output}:
\begin{enumerate}
\item \textbf{Location coefficients} ($\beta_1, \beta_2, \beta_3$): do FCI and NBFI shift the mean of future growth?
\item \textbf{Scale coefficients} ($\delta_1, \delta_2, \delta_3$): do FCI and NBFI increase the \textit{variance} of future growth?
\item \textbf{Fan chart}: predicted distribution of $y_{t+h}$ for different values of $\text{NBFI}_t$ --- showing that high NBFI leverage produces a \textit{wider} (fatter-tailed) distribution, not just a shifted one.
\item \textbf{Tail risk amplification test} (\texttt{tail\_risk\_amplification()}): quantile regression with Stress $\times$ NBFI interaction at $\tau = 0.05$, confirming that the amplification is asymmetric.

\end{enumerate}
\textbf{Key contribution}: This separates two channels: (a) NBFI leverage \textit{shifts the mean} of future growth (location effect --- already in Adrian et al. 2019); (b) NBFI leverage \textit{fattens the tails} of future growth (scale effect --- \textbf{new contribution}). If $\delta_2 > 0$ but $\beta_2 \approx 0$: NBFI leverage doesn't change average growth but makes extreme outcomes more likely.

\textbf{$\rightarrow$ Link to Module 5}: Module 4 shows NBFI leverage amplifies risk. Module 5 asks: why didn't this blow up in 2022--23? The answer lies in the network \textit{topology} --- within-OFI stress was high but didn't propagate to the real economy because banks acted as a firewall.


\bigskip\noindent\rule{\textwidth}{0.4pt}\bigskip


\subsubsection{Module 5: Resolving the Banerjee et al. (2025) Puzzle}

\textbf{Purpose}: Explain why the 2022--23 tightening cycle did \textit{not} produce a financial crisis despite high NBFI stress. Using Diebold-Yilmaz sector-level decomposition, show that within-OFI connectedness was high but OFI-to-real-economy connectedness remained low.

\textbf{Unit of observation}: System of sector-level return series $\times$ week \$t\$.

\textbf{Input variables}:


{\small
\begin{longtable}{|p{0.15\textwidth}|p{0.15\textwidth}|p{0.15\textwidth}|p{0.15\textwidth}|p{0.15\textwidth}|p{0.15\textwidth}|}
\hline
\textbf{Variable} & \textbf{Symbol} & \textbf{Definition} & \textbf{Source} & \textbf{Frequency} & \textbf{Level} \\
\hline
\endhead
OFI sector returns & $r_t^{OFI,j}$ & Returns for each OFI sub-sector \$j\$ (HFs, bond funds, MMFs, etc.) & Bloomberg / HFRX & Weekly & Sub-sector \\
\hline
Bank sector return & $r_t^{B}$ & Bank sector return & Bloomberg (KBW Bank Index) & Weekly & Sector \\
\hline
Real economy proxy & $r_t^{RE}$ & Industrial production growth, or credit growth, or non-financial corporate bond return & FRED / Bloomberg & Monthly $\rightarrow$ Weekly & Aggregate \\
\hline
Policy rate changes & $\Delta r_t^{MP}$ & Change in effective Fed funds rate & FRED & Monthly & Aggregate \\
\hline
\end{longtable}
}

\textbf{System vector}: $Y_t = (r_t^{HF}, r_t^{BondFund}, r_t^{MMF}, r_t^{IC}, r_t^{B}, r_t^{RE})'$ --- a \textbf{6-variable VAR} capturing within-OFI, OFI-to-bank, and OFI-to-real-economy links.

\textbf{Model specification} --- Standard VAR(p) with Diebold-Yilmaz GFEVD:

\[Y_t = c + A_1 Y_{t-1} + \ldots + A_p Y_{t-p} + u_t\]

From the VAR, compute the Generalised FEVD $\Theta$ at horizon \$h\$, then decompose total connectedness into:

\begin{itemize}
\item \textbf{Within-OFI connectedness}: $C^{OFI} = \frac{1}{n_{OFI}} \sum_{i,j \in OFI, i \neq j} \theta_{ij}$ --- spillovers \textit{among} OFI sub-sectors.
\item \textbf{OFI-to-bank connectedness}: $C^{OFI \to B} = \sum_{i \in OFI} \theta_{B,i}$ --- spillovers \textit{from} OFIs \textit{to} banks.
\item \textbf{OFI-to-real economy connectedness}: $C^{OFI \to RE} = \sum_{i \in OFI} \theta_{RE,i}$ --- spillovers \textit{from} OFIs \textit{to} the real economy.

\end{itemize}
\textbf{Estimation details}:
\begin{itemize}
\item Standard VAR estimated by OLS (not quantile VAR --- this module uses mean connectedness, not tail connectedness).
\item \texttt{compute\_connectedness()} in \texttt{src/analysis/systemic\_risk.py} with lags \$= 4\$, horizon \$h = 10\$.
\item \texttt{aggregate\_connectedness\_by\_sector()} aggregates the $\Theta$ matrix to sector blocks.
\item Rolling estimation (window = 60 weeks) to track how within-OFI vs. OFI-to-RE connectedness evolves over 2020--2024.

\end{itemize}
\textbf{Output}:
\begin{enumerate}
\item \textbf{Sector-block GFEVD heatmap}: shows within-OFI block is ''hot'' (high spillovers) while OFI$\rightarrow$RE block is ''cold'' (low spillovers) during 2022--23.
\item \textbf{Rolling decomposition}: time series of $C^{OFI}_t$, $C^{OFI \to B}_t$, and $C^{OFI \to RE}_t$. Expected: within-OFI spikes in 2022 but OFI$\rightarrow$RE stays flat.
\item \textbf{Narrative}: Post-Basel III bank capital buffers acted as a \textit{firewall}, absorbing OFI stress without transmitting it to the real economy. This is consistent with bank capital ratios remaining above regulatory minima throughout the tightening cycle.

\end{enumerate}
\textbf{Key insight}: The puzzle resolves because contagion has two stages: (1) within-NBFI stress propagation (high in 2022--23), and (2) NBFI-to-real-economy transmission (low in 2022--23 due to bank resilience). Standard measures that aggregate both stages miss this.


\bigskip\noindent\rule{\textwidth}{0.4pt}\bigskip


\subsubsection{Module Linkage Map (Project 2)}


\begin{adjustbox}{max width=\textwidth}
\begin{minipage}{1.5\textwidth}
\begin{verbatim}
                    +----------------------------------+
                    |  Module 1: VaR-Constraint Channel |
  Fed Z.1 leverage  |  beta_s procyclicality by sub-sector |    Output: beta_s ranking +
  data              |  + Fire-sale simulation           |    amplification ratio
                    +----------+---------+-------------+
                               |         |
              +----------------+         +----------------+
              v                                           v
+------------------------------+     +--------------------------------+
|  Module 2: DCC-GARCH         |     |  Module 3: Quantile VAR        |
|  Time-varying correlations   |     |  Tail-specific MP transmission |
|  (mean dynamics)             |     |  IRFs + connectedness          |
|  Q: Do bank-NBFI corr spike |     |  Q: Is C(0.05) >> C(0.50)?    |
|  during stress?              |     |                                |
+----------+-------------------+     +-----------+------------------+
           | Confirms regime-                     | Confirms tail
           | dependence                           | amplification
           v                                      v
+------------------------------+     +--------------------------------+
|  Module 4: Location-Scale    |     |  Module 5: Banerjee Puzzle     |
|  GaR                         |<----|  Within-OFI vs OFI→RE         |
|  Q: Does NBFI leverage       |     |  connectedness decomposition   |
|  fatten tails (variance)?    |     |  Q: Why no crisis in 2022-23? |
+------------------------------+     +--------------------------------+
\end{verbatim}
\end{minipage}
\end{adjustbox}


\subsection{Data Sources \& Sample}


{\small
\begin{longtable}{|p{0.15\textwidth}|p{0.15\textwidth}|p{0.33\textwidth}|p{0.22\textwidth}|}
\hline
\textbf{Data Source} & \textbf{Variables} & \textbf{Frequency} & \textbf{Coverage} \\
\hline
\endhead
Flow of Funds (Fed Z.1) & Sector-level assets, liabilities, leverage & Quarterly & 1980--present \\
\hline
ECB Investment Fund Statistics & EU fund AUM, flows, leverage by type & Quarterly & 2009--present \\
\hline
EPFR Global & Fund flows by type, country, asset class & Weekly/Monthly & 2000--present \\
\hline
Bloomberg / Refinitiv & Sector return indices (HF, MMF, insurance, pension) & Daily & 2000--present \\
\hline
BIS Credit Statistics & Credit to private sector, spreads & Quarterly & 1999--present \\
\hline
FRED & Fed funds rate, term spreads, GDP, CPI & Mixed & 1960--present \\
\hline
Banerjee et al. (2025) & Replication data for BIS Bulletin 116 & Quarterly & 2000--2024 \\
\hline
\end{longtable}
}

\textbf{Sample period}: 2000Q1--2024Q4 for most modules. DCC-GARCH uses daily data (2000--2024). Quantile VAR uses quarterly data.

\textbf{Sub-sector classification}: Hedge funds (HFs), money market funds (MMFs), bond mutual funds, equity mutual funds, ETFs, insurance companies (ICs), pension funds (PFs), finance companies, securitisation vehicles.


\subsection{Literature Positioning}

\begin{enumerate}
\item \textbf{Monetary policy \& financial intermediaries}: Extends Adrian \& Shin (2010, 2014) VaR-constraint channel from banks to NBFIs. Tests whether the leverage-procyclicality mechanism operates more strongly in specific NBFI sub-sectors.

\item \textbf{NBFI amplification}: Directly builds on Banerjee, Hofmann, Ng \& Pinter (2025) aggregate finding. Our sub-sector decomposition and quantile methods address their acknowledged limitation of wide confidence bands.

\item \textbf{Quantile VAR}: Adapts Ando, Greenwood-Nimmo \& Shin (2022) from cross-country to cross-sector application. First use of quantile connectedness for monetary policy transmission analysis.

\item \textbf{Location-scale models}: Extends Adrian, Boyarchenko \& Giannone (2019) GaR framework by modelling the scale (variance) channel separately, testing whether NBFI leverage fattens both tails.

\end{enumerate}

\subsection{Identification \& Robustness}

\begin{itemize}
\item \textbf{Monetary policy shock identification}: High-frequency identification (Gurkaynak, Sack \& Swanson 2005) using Fed funds futures around FOMC announcements. Robustness: Romer \& Romer (2004) narrative shocks, Jarocinski \& Karadi (2020) sign-restricted shocks.
\item \textbf{Endogeneity of NBFI leverage}: NBFI leverage is instrumented with lagged regulatory capital ratios of connected banks (supply-side shifters).
\item \textbf{Quantile VAR lag selection}: BIC-optimal lag length at each quantile; robustness across $p \in \{1, 2, 4\}$.
\item \textbf{Small-sample inference}: Bootstrap confidence intervals (1000 replications) for quantile IRFs and connectedness measures.

\end{itemize}

\subsection{Expected Contributions}

\begin{enumerate}
\item Sub-sector decomposition: hedge funds drive amplification, pension funds dampen it.
\item Quantile connectedness reveals tail-specific MP amplification invisible to linear methods.
\item Location-scale evidence: NBFI leverage fattens both tails (variance effect), not just shifts the mean.
\item Resolution of the 2022--23 resilience puzzle via network topology analysis.

\end{enumerate}

\bigskip\noindent\rule{\textwidth}{0.4pt}\bigskip


\section{Project 3: Contagion Across Borders}

\textbf{Subtitle:} How NBFI Stress in One Jurisdiction Spills Over via the Global Financial Cycle


\subsection{Motivation}

Cross-border NBFI linkages create a transmission mechanism for financial stress
through three interconnected channels: (1) dollar funding, where NBFIs borrow
in USD via repo and FX swaps; (2) portfolio rebalancing by global investment
funds whose procyclical flows amplify local shocks; and (3) the global financial
cycle (Rey 2013), a common factor in risky asset prices that NBFIs both respond
to and amplify.


\subsection{Research Questions}

\textbf{RQ1.} Is cross-border NBFI connectedness asymmetric --- dramatically higher in
the left tail (stress) than at the median (normal times)?

\textbf{RQ2.} Does dollar funding stress transmit cross-border, and is this
transmission nonlinear (stronger in the tails)?

\textbf{RQ3.} Are NBFI portfolio flows more volatile and more cross-border-connected
than bank flows?

\textbf{RQ4.} Does NBFI penetration amplify the transmission of the global financial
cycle to local financial conditions? Is this robust to IV estimation?

\textbf{RQ5.} Can a hedge fund failure cascade through prime brokerage links to
generate cross-border deleveraging?


\subsection{Data Sources \& Sample}


{\small
\begin{longtable}{|p{0.15\textwidth}|p{0.15\textwidth}|p{0.33\textwidth}|p{0.22\textwidth}|}
\hline
\textbf{Data Source} & \textbf{Variables} & \textbf{Frequency} & \textbf{Coverage} \\
\hline
\endhead
BIS Locational Banking Stats & Cross-border claims by sector \& country & Quarterly & 2000--present \\
\hline
EPFR Global & Cross-border fund flows by type \& destination & Monthly & 2005--present \\
\hline
IMF CPIS & Portfolio investment positions by country pair & Annual & 2001--present \\
\hline
BIS OTC Derivatives & FX swap/forward outstanding by currency & Semiannual & 2004--present \\
\hline
Bloomberg & CIP basis (cross-currency basis swaps), VIX, equity indices & Daily & 2000--present \\
\hline
Miranda-Agrippino \& Rey (2020) & Global Financial Cycle factor (updated) & Monthly & 1990--present \\
\hline
FSB Global Monitoring & NBFI penetration ratio by country & Annual & 2002--present \\
\hline
\end{longtable}
}

\textbf{Sample period}: 2005Q1--2024Q4 for cross-country quantile connectedness (EPFR flow data availability). Extended to 2000 for GFC panel regressions using BIS data.

\textbf{Country coverage}: G20 + selected financial centres (US, UK, EA, JP, CH, AU, CA, KR, SG, HK, BR, MX, IN, ZA, TR, RU, CN). Separate EME sub-sample for GFC amplification tests.


\subsection{Literature Positioning}

\begin{enumerate}
\item \textbf{Global financial cycle}: Extends Rey (2013) and Miranda-Agrippino \& Rey (2020) by testing whether NBFI penetration amplifies GFC transmission. First causal evidence using IV (US monetary policy shocks $\rightarrow$ GFC $\rightarrow$ local conditions, with NBFI interaction).

\item \textbf{Cross-border contagion}: Complements Forbes \& Warnock (2012) on capital flow surges/stops and Broner, Didier, Erce \& Schmukler (2013) on gross flows. Adds NBFI-specific decomposition and quantile methods.

\item \textbf{Dollar funding}: Builds on Avdjiev, Du, Koch \& Shin (2019) and Eguren-Martin, Ossandon Busch \& Reinhardt (2024) on dollar funding channels. Tests nonlinear transmission via quantile regression.

\item \textbf{Network contagion}: Agent-based simulation follows Eisenberg \& Noe (2001) clearing mechanism with Kyle (1985) price impact, extending Cont \& Schaanning (2017) fire-sale framework to cross-border setting.

\end{enumerate}

\subsection{Identification \& Robustness}

\begin{itemize}
\item \textbf{IV for GFC factor}: US monetary policy shocks (Jarocinski \& Karadi 2020) instrument the GFC, addressing reverse causality from local conditions to global factor.
\item \textbf{NBFI penetration endogeneity}: NBFI share instrumented with (a) legal origin (La Porta et al. 1998) for cross-section, (b) lagged pension reform dummies for within-country variation.
\item \textbf{Alternative connectedness measures}: Barunik \& Krehlik (2018) frequency-domain connectedness as robustness check on time-domain results.
\item \textbf{EME vs AE heterogeneity}: Full interaction models allowing different coefficients for advanced and emerging economies.

\end{itemize}

\subsection{Empirical Strategy}

Project 3 traces three contagion channels --- each with a specific NBFI mechanism --- then unifies them in a cascade simulation:


\begin{adjustbox}{max width=\textwidth}
\begin{minipage}{1.5\textwidth}
\begin{verbatim}
Module 1 (Dollar Funding / FX Hedge)  --→  CIP stress amplified by NBFI hedging demand
Module 2 (Procyclical Rebalancing)    --→  Fire-sale flows from open-end fund redemptions
Module 3 (Cross-Border Amplification) --→  Channel-identified panel: how NBFI integration transmits GFC/GTC
Module 4 (Cascade Simulation)         --→  Agent-based contagion: HF failure → PB → cross-border losses
\end{verbatim}
\end{minipage}
\end{adjustbox}


\bigskip\noindent\rule{\textwidth}{0.4pt}\bigskip


\subsubsection{Module 1: Dollar Funding Channel}

\textbf{Purpose}: Test whether dollar funding stress (measured by the CIP basis) is amplified in countries whose financial sectors have larger structural dollar hedging needs, and whether this amplification is nonlinear (stronger in the tails).


\paragraph{Economic Logic}

The CIP basis is determined by the balance between supply and demand for dollar hedging in the FX swap market. What matters is not generic ''NBFI size'' but specifically \textbf{how much of country \$c\$'s financial sector relies on FX-hedged dollar assets or dollar funding via swaps}:

\begin{itemize}
\item \textbf{Pension funds and insurance companies} in countries like Japan, the Netherlands, and Korea hold large portfolios of US Treasuries and corporate bonds, hedged via 3-month FX forwards that must be rolled quarterly. This creates structural, recurring demand for dollar hedging.
\item \textbf{NBFI dollar borrowing}: Institutions that use FX swaps as synthetic dollar funding (borrow dollars, lend local currency) are directly exposed when the swap market seizes up.

\end{itemize}
When stress hits (VIX spikes, equity markets fall), these institutions face margin calls on their FX forwards and compete for scarce dollar liquidity. Countries with larger structural dollar hedging demand see their CIP basis blow out more --- not because NBFIs are ''big'' but because their hedging activity creates concentrated demand that overwhelms dealer capacity during stress.

\textbf{Unit of observation}: Currency pair \$c\$ $\times$ month \$t\$.

\textbf{Input variables}:


{\small
\begin{longtable}{|p{0.15\textwidth}|p{0.15\textwidth}|p{0.15\textwidth}|p{0.15\textwidth}|p{0.15\textwidth}|p{0.15\textwidth}|}
\hline
\textbf{Variable} & \textbf{Symbol} & \textbf{Definition} & \textbf{Source} & \textbf{Frequency} & \textbf{Level} \\
\hline
\endhead
CIP basis & $\text{CIP}_{c,t}$ & Cross-currency basis swap spread for currency \$c\$ vs USD (3-month tenor). Negative = dollar funding premium & Bloomberg & Daily $\rightarrow$ Monthly avg & Currency pair \\
\hline
FX hedge demand & $\text{FXHedge}_{c,t}$ & Structural dollar hedging need of country \$c\$'s financial sector (see below) & BIS Triennial Survey; IMF CPIS; national sources & Semi-annual / Annual & Country \\
\hline
VIX & $\text{VIX}_t$ & CBOE Volatility Index & FRED & Daily $\rightarrow$ Monthly avg & Global \\
\hline
USD broad index & --- & Trade-weighted USD index & FRED (DTWEXBGS) & Daily $\rightarrow$ Monthly avg & Global \\
\hline
Local equity return & $r_{c,t}^{EQ}$ & Country \$c\$ equity index return & Bloomberg & Daily $\rightarrow$ Monthly & Country \\
\hline
\end{longtable}
}

\textbf{Measuring $\text{FXHedge}_{c,t}$} --- three alternatives (from most to least precise):
\begin{enumerate}
\item \textbf{Outstanding FX swap/forward positions} of country \$c\$'s non-bank financial sector (from BIS Triennial Survey --- available for major currencies, semi-annual)
\item \textbf{Cross-border bond holdings that need hedging}: country \$c\$'s institutional investors' holdings of foreign-currency bonds (from IMF Coordinated Portfolio Investment Survey), scaled by estimated hedge ratios
\item \textbf{Fallback}: NBFI foreign assets as \% of GDP from FSB Global Monitoring Report (coarser but broader country coverage)

\end{enumerate}
\textbf{Model specification} --- Quantile regression:

\[Q_\tau(\text{CIP}_{c,t}) = \alpha_c + \beta_1(\tau) \cdot \text{VIX}_t + \beta_2(\tau) \cdot r_{c,t}^{EQ} + \beta_3(\tau) \cdot \text{FXHedge}_{c,t} + \beta_4(\tau) \cdot (\text{VIX}_t \times \text{FXHedge}_{c,t}) + \varepsilon_{c,t}\]

\textbf{Key test}: $\beta_4(\tau = 0.05) \neq 0$ but $\beta_4(\tau = 0.50) \approx 0$. Countries with larger structural dollar hedging needs suffer larger CIP blowouts during stress --- but not during normal times. The mechanism: stress triggers a rush to roll hedges, demand overwhelms dealer capacity, and the basis blows out in proportion to pre-existing hedging demand.

\textbf{Estimation}: \texttt{statsmodels.QuantReg} at $\tau \in \{0.05, 0.25, 0.50, 0.75, 0.95\}$, with country fixed effects. Standard errors clustered by country.

\textbf{Output}: Coefficient table across quantiles; scatter plot of CIP basis vs. VIX for high- vs. low-FXHedge countries.


\bigskip\noindent\rule{\textwidth}{0.4pt}\bigskip


\subsubsection{Module 2: Procyclical Rebalancing Channel}

\textbf{Purpose}: Test whether cross-border portfolio flows are amplified by the share of assets managed by open-ended investment funds (OEFs), which face redemption-driven forced selling.


\paragraph{Economic Logic}

Portfolio flows net to zero in aggregate --- every sale is someone else's purchase. The amplification question is whether the \textit{composition} of intermediation matters. Open-ended funds face a structural vulnerability: when investors redeem shares, the fund must liquidate assets regardless of price. This creates a procyclical channel:

\begin{enumerate}
\item Global shock (VIX spike, GFC/GTC deterioration) $\rightarrow$ investor redemptions from OEFs
\item OEFs forced to sell cross-border assets $\rightarrow$ outflows from recipient country \$c\$
\item Fire-sale prices overshoot $\rightarrow$ further NAV decline $\rightarrow$ more redemptions (feedback loop)
\item Countries where OEFs intermediate a larger share of cross-border flows see larger amplification

\end{enumerate}
By contrast, closed-ended funds, pension funds, and banks do not face the same redemption pressure --- they can ride out mark-to-market losses.


\paragraph{Estimation}

\textbf{Test 1 --- Flow amplification}:

\[\text{Flow}_{c,t}^{\text{NBFI}} = \alpha_c + \beta_1 \cdot \text{GFC}_t + \beta_2 \cdot \text{OEF\_Share}_{c} + \beta_3 \cdot (\text{GFC}_t \times \text{OEF\_Share}_{c}) + \gamma X_{c,t} + \varepsilon_{c,t}\]

where $\text{Flow}_{c,t}^{\text{NBFI}}$ = net cross-border portfolio flows intermediated by NBFIs (from BIS LBS), $\text{OEF\_Share}_{c}$ = fraction of country \$c\$'s cross-border portfolio assets managed by open-ended funds (from FSB NBFI monitoring data), and $X_{c,t}$ = controls (GDP growth, exchange rate, interest rate differential).

\textbf{Key test}: $\beta_3 < 0$ --- during global stress ($\text{GFC}_t$ high), countries with higher OEF share see larger outflows.

\textbf{Test 2 --- Differential behavior, NBFI vs. bank flows}:

Run the same specification separately for $\text{Flow}_{c,t}^{\text{NBFI}}$ and $\text{Flow}_{c,t}^{\text{Bank}}$ (bank-intermediated cross-border flows). The amplification coefficient $\beta_3$ should be significant only for NBFI flows.

\textbf{Test 3 --- Price impact}: Regress asset price changes on the predicted component of NBFI outflows (instrumented by $\text{GFC}_t \times \text{OEF\_Share}_{c}$) to confirm that fire-sale flows move prices.

\textbf{Data}: BIS Locational Banking Statistics (cross-border flows by counterparty type), FSB Global Monitoring Report (OEF shares by country), EPFR fund-level flows (robustness).

\textbf{Output}: Coefficient table for $\beta_3$ across flow types; scatter plot of predicted NBFI outflows vs. asset price declines during stress episodes.


\bigskip\noindent\rule{\textwidth}{0.4pt}\bigskip


\subsubsection{Module 3: Cross-Border Amplification Panel}

\textbf{Purpose}: Identify \textit{which NBFI integration channels} amplify the transmission of global shocks to domestic financial conditions, using the counterparty decomposition developed in Project 1.


\paragraph{Economic Logic}

The key insight from the user's critique is that NBFI ''penetration'' is not a single number. The effects of shocks depend on \textit{how} NBFIs are integrated in the financial sector:

\begin{itemize}
\item \textbf{NBFI borrows from foreign banks} (for-borrow): creates dollar funding vulnerability (links to Module 1)
\item \textbf{NBFI borrows from domestic banks} (bank-borrow): creates domestic credit risk transmission
\item \textbf{NBFI lends to private sector} (priv-lend): creates real economy exposure through non-bank credit
\item \textbf{NBFI participates in FX hedging} (fx-hedge): creates FX swap market stress channel (links to Project 4)

\end{itemize}
Each channel has a different amplification mechanism, and lumping them together in a single ''NBFI size'' variable would obscure the results.


\paragraph{Estimation}

\textbf{Main specification} --- Panel regression with NBFI counterparty decomposition:

\[y_{c,t} = \alpha_c + \delta_t + \sum_{k} \beta_k \cdot \text{NBFI}_{c,t-1}^{(k)} + \sum_{k} \gamma_k \cdot (\text{Shock}_t \times \text{NBFI}_{c,t-1}^{(k)}) + \theta X_{c,t} + \varepsilon_{c,t}\]

where:
\begin{itemize}
\item $y_{c,t}$ = domestic financial conditions index (credit spread, or Growth-at-Risk 5th percentile from quantile regression of GDP growth on financial conditions)
\item $\text{NBFI}_{c,t-1}^{(k)}$ = lagged NBFI counterparty shares, $k \in \{$for-borrow, bank-borrow, priv-lend, fx-hedge$\}$
\item $\text{Shock}_t$ = global shock measure. Two alternatives:
\item GFC (Global Financial Cycle): first principal component of risky asset prices (Rey 2015)
\item GTC (Global Tail Cycle): common factor in left-tail co-movements (extracted from the QVAR quantile factor model in Project 2, Module 3)
\item $\delta_t$ = time fixed effects, which \textbf{absorb the direct effect of GFC/GTC} --- identification comes purely from \textit{cross-country variation} in NBFI structure
\item $X_{c,t}$ = controls: GDP growth, policy rate, exchange rate regime, capital account openness

\end{itemize}
\textbf{Key tests}:
\begin{itemize}
\item $\gamma_{\text{for-borrow}} < 0$: countries where NBFIs borrow more from foreign banks see worse financial conditions during global stress
\item $\gamma_{\text{priv-lend}} < 0$: countries where NBFIs intermediate more credit to the private sector see larger real economy impact
\item Compare magnitudes: which channel dominates?

\end{itemize}
\textbf{Instrumental variable strategy}: To address endogeneity of NBFI structure, instrument $\text{Shock}_t$ with US monetary policy surprises (Jarocinski \& Karadi 2020). The exclusion restriction: US MP shocks affect country \$c\$ only through global financial conditions, not directly through bilateral NBFI linkages.

\textbf{Extensions}:
\begin{itemize}
\item Replace $y_{c,t}$ with 5th-percentile GDP growth (Growth-at-Risk) to test tail effects
\item Interact with state-dependence indicator: $\mathbb{1}[\text{VIX}_t > \text{VIX}_{80\text{th pctile}}]$ to test whether amplification is non-linear
\item Use GTC instead of GFC as the shock measure --- if GTC interactions are stronger, it confirms that tail co-movements (not just average cycles) drive amplification

\end{itemize}
\textbf{Data}: BIS LBS (NBFI counterparty decomposition), IMF IFS (financial conditions), Jarocinski-Karadi MP shocks, GFC from Project 2 (Module 2 DCC-GARCH or PCA of risky asset returns), GTC from Project 2 (Module 3 QVAR quantile factor extraction).

\textbf{Output}: Coefficient table for $\gamma_k$ across channels; bar chart comparing amplification magnitudes by NBFI integration type; GFC vs. GTC comparison.


\bigskip\noindent\rule{\textwidth}{0.4pt}\bigskip


\subsubsection{Module 4: Contagion Cascade Simulation}

\textbf{Purpose}: Agent-based simulation of a cross-border contagion event: US hedge fund failure $\rightarrow$ prime broker losses $\rightarrow$ multi-round forced deleveraging $\rightarrow$ cross-border loss propagation.

\textbf{Unit of observation}: Simulated agents (hedge funds, prime brokers, banks across countries).

\textbf{Input variables (calibration)}:


{\small
\begin{longtable}{|p{0.22\textwidth}|p{0.33\textwidth}|p{0.33\textwidth}|}
\hline
\textbf{Parameter} & \textbf{Value} & \textbf{Source} \\
\hline
\endhead
US hedge fund sector & 20 funds, NAV \$2M total, leverage 3--15$\times$ & Calibrated from SEC Form PF \\
\hline
Prime brokers & GS, JPM, MS, CS, DB & Major PB list \\
\hline
Market depth & \\$300K & Calibrated to average equity market depth \\
\hline
Margin requirement & 10\% of gross exposure & Industry standard \\
\hline
Margin procyclicality & 0.5 (margin rises 50\% of price drop) & Brunnermeier \& Pedersen (2009) \\
\hline
Initial shock & -4\% across all assets & Scenario \\
\hline
\end{longtable}
}

\textbf{Simulation mechanism} (Eisenberg-Noe 2001 + Kyle 1985):
\begin{enumerate}
\item Exogenous shock to hedge fund portfolios (-4\%).
\item Funds with insufficient margin receive margin calls from prime brokers.
\item Funds that cannot meet margin are forced to deleverage: sell assets proportional to excess leverage.
\item Aggregate selling $\rightarrow$ price impact $\rightarrow$ further losses $\rightarrow$ more margin calls.
\item Funds that default impose credit losses on their prime brokers.
\item Prime brokers (which are cross-border banks) transmit losses to other jurisdictions.
\item Iterate until convergence or max rounds.

\end{enumerate}
\textbf{Estimation details}: Implemented via \texttt{PrimeBrokerageContagion} class in \texttt{src/models/nbfi\_subsectors.py}. \texttt{build\_hedge\_fund\_sector()} creates the fund universe. \texttt{.run(initial\_shock=-0.04)} executes the cascade.

\textbf{Output}:
\begin{enumerate}
\item Round-by-round dynamics: market price, total NAV, gross exposure, average leverage, fund defaults, PB losses.
\item Cross-border decomposition: which prime brokers (and hence which countries) absorb the most losses.
\item 3-panel chart: NAV path, leverage path, default count.

\end{enumerate}

\bigskip\noindent\rule{\textwidth}{0.4pt}\bigskip


\subsubsection{Module Linkage Map (Project 3)}


\begin{adjustbox}{max width=\textwidth}
\begin{minipage}{1.5\textwidth}
\begin{verbatim}
+-------------------------------+
| Module 1: Dollar Funding      | FXHedge x VIX → CIP stress at tails
| (country-level quantile reg.) |
+---------------+---------------+
                | feeds fx-hedge channel into Module 3
                |
+---------------+---------------+
| Module 2: Procyclical         | OEF_Share x GFC → fire-sale outflows
| Rebalancing (flow regressions)|
+---------------+---------------+
                | predicted flows enter Module 3 as instrument
                |
+---------------+---------------+
| Module 3: Cross-Border        | NBFI counterparty decomposition x Shock
| Amplification Panel           | → domestic financial conditions / GaR
| (channel-identified)          |
+---------------+---------------+
                | calibrates shock amplification for simulation
                |
+---------------+---------------+
| Module 4: Cascade Simulation  | HF failure → PB losses → cross-border
| (Eisenberg-Noe + Kyle)        | deleveraging with amplification from M1-M3
+-------------------------------+
\end{verbatim}
\end{minipage}
\end{adjustbox}


\subsection{Expected Contributions}

\begin{enumerate}
\item Left-tail connectedness across borders is 40--60\% higher than median connectedness.
\item NBFI portfolio flows are systematically more volatile and more cross-border-connected than bank flows.
\item NBFI penetration amplifies GFC transmission to local credit conditions (robust to IV with US MP shocks).
\item Location-scale evidence: NBFI flows fatten the tails of EME return distributions.

\end{enumerate}

\bigskip\noindent\rule{\textwidth}{0.4pt}\bigskip


\section{Project 4: FX Hedging as a Contagion Channel}

\textbf{Subtitle:} How NBFI Currency Risk Management Transmits Global Financial Shocks


\subsection{Motivation}

The global FX derivatives market has grown to \\$75 trillion in outstanding
notional (BIS OTC statistics, end-2024), driven primarily by NBFIs hedging
cross-border bond investments. Combining Rey, Stavrakeva \& Tang (2024) on
currency centrality with Nenova, Schrimpf \& Shin (2025) on FX derivatives and
NBFI hedging, we identify a complete contagion channel:

**Equity shocks $\rightarrow$ FX movements $\rightarrow$ hedging cost changes $\rightarrow$ NBFI portfolio
adjustment $\rightarrow$ cross-border bond flow reversal $\rightarrow$ feedback to asset prices**

Critically, yield curve movements partially offset FX-induced hedging cost
changes in normal times (self-stabilising), but during stress all forces
reinforce --- creating nonlinear amplification detectable only through quantile
methods.


\subsection{Data Sources \& Sample}


{\small
\begin{longtable}{|p{0.15\textwidth}|p{0.15\textwidth}|p{0.33\textwidth}|p{0.22\textwidth}|}
\hline
\textbf{Data Source} & \textbf{Variables} & \textbf{Frequency} & \textbf{Coverage} \\
\hline
\endhead
BIS OTC Derivatives Stats & FX swap/forward outstanding by currency \& sector & Semiannual & 2004--present \\
\hline
Bloomberg & Cross-currency basis swaps (3M, 1Y, 5Y), equity indices, VIX & Daily & 2000--present \\
\hline
Refinitiv & Government bond yields (2Y, 5Y, 10Y) for 20+ countries & Daily & 2000--present \\
\hline
EPFR Global & Cross-border bond fund flows by country & Monthly & 2005--present \\
\hline
IMF CPIS & Cross-border bond holdings by country pair & Annual & 2001--present \\
\hline
Rey, Stavrakeva \& Tang (2024) & Currency centrality measures, replication data & Monthly & 1999--2023 \\
\hline
Nenova, Schrimpf \& Shin (2025) & FX derivatives and NBFI hedging data & Quarterly & 2010--2024 \\
\hline
\end{longtable}
}

\textbf{Sample period}: 2005M1--2024M12 for quantile regressions (cross-currency basis swap availability). Semiannual panel for BIS derivatives data.

\textbf{Currency pairs}: USD/EUR, USD/JPY, USD/GBP, USD/AUD, USD/CAD, USD/CHF (G10 majors) + USD/KRW, USD/MXN, USD/BRL, USD/ZAR (key EME pairs with liquid CIP basis data).


\subsection{Literature Positioning}

\begin{enumerate}
\item \textbf{FX hedging and financial stability}: Directly integrates two recent contributions --- Rey, Stavrakeva \& Tang (2024) on equity-FX transmission and currency centrality, and Nenova, Schrimpf \& Shin (2025) on FX derivatives and NBFI hedging demand. The ''offsetting forces'' hypothesis is our key theoretical contribution.

\item \textbf{CIP deviations}: Extends Du, Tepper \& Verdelhan (2018) and Avdjiev, Du, Koch \& Shin (2019) on post-GFC CIP violations. We show that CIP deviations are not just a funding cost anomaly but a systemic risk transmission channel during stress.

\item \textbf{Delta-CoVaR}: Applies Adrian \& Brunnermeier (2016) to the FX hedging channel specifically. First application of CoVaR to measure systemic contribution of a specific transmission mechanism rather than an institution.

\item \textbf{Maturity mismatch}: Builds on Brunnermeier, Nagel \& Pedersen (2009) carry trade analysis. NBFI hedging mismatch (3-month FX swaps rolled to hedge 10-year bond positions) creates systemic rollover risk analogous to bank maturity mismatch.

\end{enumerate}

\subsection{Identification \& Robustness}

\begin{itemize}
\item \textbf{Offsetting forces test}: Formal Wald test of $H_0: \beta_4(\tau=0.05) = \beta_4(\tau=0.50)$ using Koenker \& Bassett (1982) framework. Bootstrap p-values for interquantile differences.
\item \textbf{Endogeneity of CIP basis}: CIP basis instrumented with central bank swap line announcements (exogenous supply-side shocks to FX funding).
\item \textbf{Hedging demand proxy}: NBFI cross-border bond holdings (IMF CPIS) $\times$ average hedge ratio (BIS survey) as measure of hedging demand.
\item \textbf{Alternative stress measures}: Results robust to replacing VIX with MOVE (bond volatility), TED spread, or financial conditions indices.

\end{itemize}

\subsection{Research Questions}

\textbf{RQ1.} Do the offsetting forces between equity-driven FX movements and yield
curve dynamics break down in the tails of the distribution? (Quantile regression
of CIP basis.)

\textbf{RQ2.} Does FX hedging stress have systemic risk implications for cross-border
bond flows? (Delta-CoVaR framework.)

\textbf{RQ3.} Does the equity-FX-CIP-bonds transmission chain activate asymmetrically
during crises versus normal times? (Quantile connectedness: $\tau = 0.05$ vs $\tau = 0.50$.)

\textbf{RQ4.} Do countries with higher NBFI penetration experience stronger contagion
through the FX hedging channel? (Panel regression with NBFI interaction + IV.)


\bigskip\noindent\rule{\textwidth}{0.4pt}\bigskip


\section{Timeline \& Deliverables}


{\small
\begin{longtable}{|p{0.15\textwidth}|p{0.15\textwidth}|p{0.33\textwidth}|p{0.22\textwidth}|}
\hline
\textbf{Phase} & \textbf{Period} & \textbf{Projects} & \textbf{Deliverables} \\
\hline
\endhead
\textbf{Phase 1} & Q1--Q2 2026 & P1 + P2 (parallel) & Working papers; P1 implied leverage dataset; P2 quantile connectedness estimates \\
\hline
\textbf{Phase 2} & Q3--Q4 2026 & P3 + P4 (parallel) & Working papers; P3 cross-country connectedness database; P4 offsetting forces evidence \\
\hline
\textbf{Phase 3} & Q1 2027 & Integration & Synthesis paper linking all four projects; unified policy brief \\
\hline
\textbf{Ongoing} & Throughout & All & Conference presentations, seminar feedback, revisions \\
\hline
\end{longtable}
}

\textbf{Sequencing rationale}: P1 and P2 can proceed in parallel (different data, complementary methods). P3 and P4 benefit from P1/P2 outputs (leverage estimates, sub-sector results) but can begin data assembly immediately. The integration phase produces the synthesis that ties all four dimensions together.


\bigskip\noindent\rule{\textwidth}{0.4pt}\bigskip


\section{Policy Implications}

The research programme speaks directly to ongoing regulatory debates:

\begin{enumerate}
\item \textbf{NBFI leverage monitoring (P1)}: The implied leverage measure and HLI provide regulators with a prototype monitoring tool that complements existing FSB surveillance. Could inform the design of NBFI leverage reporting requirements currently under discussion (FSB 2023 NBFI roadmap).

\item \textbf{Monetary policy \& financial stability (P2)}: Evidence that NBFI leverage amplifies MP transmission --- particularly in the tails --- has implications for how central banks calibrate tightening cycles. The 2022--23 resilience finding suggests post-GFC bank regulation worked but created new vulnerabilities in the NBFI sector.

\item \textbf{Cross-border macroprudential coordination (P3)}: If NBFI penetration amplifies GFC transmission, countries with large NBFI sectors face stronger exposure to external financial shocks. Supports arguments for cross-border macroprudential coordination (e.g., reciprocity of countercyclical capital buffers for NBFI exposures).

\item \textbf{FX derivatives regulation (P4)}: The maturity mismatch in NBFI FX hedging (3-month swaps for 10-year bonds) creates systemic rollover risk. Supports proposals for longer-tenor hedging requirements or margin buffers calibrated to stress scenarios.

\item \textbf{Data gaps (all projects)}: All four projects highlight the inadequacy of current NBFI data. Supports FSB/BIS initiatives for enhanced NBFI reporting, particularly on leverage, counterparty exposures, and cross-border linkages.

\end{enumerate}

\bigskip\noindent\rule{\textwidth}{0.4pt}\bigskip


\section{References (Selected)}

\begin{itemize}
\item Acemoglu, D., Ozdaglar, A., \& Tahbaz-Salehi, A. (2015). Systemic risk and stability in financial networks. \textit{American Economic Review}, 105(2), 564--608.
\item Adrian, T., \& Brunnermeier, M. K. (2016). CoVaR. \textit{American Economic Review}, 106(7), 1705--1741.
\item Adrian, T., Boyarchenko, N., \& Giannone, D. (2019). Vulnerable growth. \textit{American Economic Review}, 109(4), 1263--1289.
\item Adrian, T., \& Shin, H. S. (2010). Liquidity and leverage. \textit{Journal of Financial Intermediation}, 19(3), 418--437.
\item Ando, T., Greenwood-Nimmo, M., \& Shin, Y. (2022). Quantile connectedness: Modeling tail behavior in the topology of financial networks. \textit{Management Science}, 68(4), 2401--2431.
\item Avdjiev, S., Du, W., Koch, C., \& Shin, H. S. (2019). The dollar, bank leverage, and deviations from covered interest parity. \textit{American Economic Review: Insights}, 1(2), 193--208.
\item Banerjee, R., Hofmann, B., Ng, A., \& Pinter, J. (2025). Non-bank financial intermediaries and financial stability. \textit{BIS Bulletin}, 116.
\item Du, W., Tepper, A., \& Verdelhan, A. (2018). Deviations from covered interest rate parity. \textit{Journal of Finance}, 73(3), 915--957.
\item Eisenberg, L., \& Noe, T. H. (2001). Systemic risk in financial systems. \textit{Management Science}, 47(2), 236--249.
\item Forbes, K. J., \& Warnock, F. E. (2012). Capital flow waves: Surges, stops, flight, and retrenchment. \textit{Journal of International Economics}, 88(2), 235--251.
\item Greenwood, R., Landier, A., \& Thesmar, D. (2015). Vulnerable banks. \textit{Journal of Financial Economics}, 115(3), 471--485.
\item Miranda-Agrippino, S., \& Rey, H. (2020). US monetary policy and the global financial cycle. \textit{Review of Economic Studies}, 87(6), 2754--2776.
\item Nenova, T., Schrimpf, A., \& Shin, H. S. (2025). FX derivatives and NBFI hedging. \textit{BIS Working Paper}.
\item Rey, H. (2013). Dilemma not trilemma: The global financial cycle and monetary policy independence. \textit{Jackson Hole Symposium}.
\item Rey, H., Stavrakeva, V., \& Tang, J. (2024). Currency centrality and the exchange rate channel. \textit{Working Paper}.

\end{itemize}

\subsection{Empirical Strategy}

Project 4 has \textbf{three analytical layers} plus a \textbf{cascade simulation} and an \textbf{NBFI amplification panel}. The three layers test the ''offsetting forces'' hypothesis at increasing levels of complexity.


\begin{adjustbox}{max width=\textwidth}
\begin{minipage}{1.5\textwidth}
\begin{verbatim}
Layer 1 (Quantile Regression)  --→  Layer 2 (Delta-CoVaR)  --→  Layer 3 (Quantile Connectedness)  --→  NBFI Panel  --→  Cascade Sim
         ↓                                  ↓                              ↓                              ↓                  ↓
   Offsetting forces:              Systemic risk of           Full transmission         Countries with        Multi-round
   interaction breaks              FX hedging stress          chain dormant at          higher NBFI           equity → FX →
   down in tails                   on bond flows              median, active            suffer more           CIP → bonds
                                                              at tau=0.05                                      feedback loop
\end{verbatim}
\end{minipage}
\end{adjustbox}


\bigskip\noindent\rule{\textwidth}{0.4pt}\bigskip


\subsubsection{The Offsetting Forces Hypothesis}

Before detailing each layer, the key economic mechanism. The relevant shock is a \textbf{global risk-off episode} --- a broad-based flight from risky assets worldwide, typically proxied by a large S\&P 500 decline. The distinction matters: a US-specific correction (e.g., tech sector sell-off) could trigger portfolio rebalancing \textit{away} from USD. But a global risk-off event activates the dollar's safe-haven role, which is the setting for the contagion channel below.


{\small
\begin{longtable}{|p{0.22\textwidth}|p{0.33\textwidth}|p{0.33\textwidth}|}
\hline
\textbf{Force} & \textbf{Normal Times (moderate risk-off)} & \textbf{Stress (severe global risk-off)} \\
\hline
\endhead
Global risk-off $\rightarrow$ FX (Rey et al. 2024) & USD appreciates moderately (safe-haven demand outweighs portfolio rebalancing) & USD appreciates violently (dollar funding squeeze + flight to safety) \\
\hline
Yield curve & Bull steepens (front end drops on easing expectations; long end holds) $\rightarrow$ attractive hedged carry & Bull flattens/inverts (flight to quality crushes long end) $\rightarrow$ carry disappears \\
\hline
CIP deviations & Small $\rightarrow$ manageable hedging cost & Blow out $\rightarrow$ prohibitive cost \\
\hline
\textbf{Net effect} & \textbf{Partially offsetting $\rightarrow$ self-stabilising} & \textbf{All reinforcing $\rightarrow$ amplification} \\
\hline
\end{longtable}
}


\paragraph{How the Mechanism Works: A Step-by-Step Example}

Consider a European pension fund holding US Treasury bonds, hedged using 3-month FX forwards (rolled quarterly). The cost of hedging is reflected in the CIP basis.

\textbf{The shock}: A global risk-off event triggers broad equity declines. US equity markets fall sharply alongside global equities.

\textbf{Channel 1 --- Exchange rate (Rey et al. 2024)}: In a global risk-off, two forces compete. Portfolio rebalancing pushes investors to sell US equities and repatriate capital $\rightarrow$ USD selling pressure. But the safe-haven channel pushes in the opposite direction: global investors flee risky assets \textit{everywhere} (EM equities, corporate bonds, commodities) and rush into USD-denominated safe assets; simultaneously, global banks scramble to roll USD funding, creating a dollar shortage (Avdjiev et al. 2019). Empirically, the safe-haven / funding channel dominates --- the USD appreciates. EUR/USD falls. For the pension fund:
\begin{itemize}
\item Its FX hedge (short USD forward) suffers a mark-to-market loss
\item Hedging \textit{future} USD exposure becomes more expensive because the CIP basis widens as dollar demand surges

\end{itemize}
\textbf{Channel 2 --- Yield curve (Nenova et al. 2025)}: The equity shock affects monetary policy expectations, but the yield curve response differs by severity:
\begin{itemize}
\item \textit{Normal times} (moderate correction): Markets price in modest rate cuts $\rightarrow$ the \textbf{short end} (2Y) drops quickly (directly tied to policy expectations), while the \textbf{long end} (10Y) is anchored by term premia and long-run growth expectations $\rightarrow$ 10Y-2Y spread \textbf{widens} (bull steepening) $\rightarrow$ hedged carry (long-end yield minus short-term hedging cost) remains \textbf{attractive} $\rightarrow$ pension fund keeps hedging and buying US bonds
\item \textit{Stress} (severe crash): Flight to quality drives aggressive buying of long-term Treasuries $\rightarrow$ the \textbf{long end} drops sharply $\rightarrow$ 10Y-2Y spread \textbf{narrows or inverts} (bull flattening) $\rightarrow$ hedged carry \textbf{collapses} $\rightarrow$ the business case for holding hedged US bonds evaporates

\end{itemize}
\textbf{Channel 3 --- CIP deviations}: The CIP basis (FX hedging cost above covered interest parity) is normally small (5--15 bps). In stress:
\begin{itemize}
\item Banks that intermediate FX swaps pull back (their own VaR constraints bind)
\item Demand for USD hedging surges (everyone hedges in the same direction)
\item The basis \textbf{blows out} (50--100+ bps) $\rightarrow$ hedging becomes prohibitively expensive

\end{itemize}
\textbf{The key asymmetry}:

\textit{Normal times} (say, a 5\% global equity correction):
\begin{itemize}
\item USD appreciates modestly (safe-haven demand slightly exceeds portfolio rebalancing) $\rightarrow$ hedging cost up slightly (\textbf{negative})
\item Yield curve bull-steepens (short end drops, long end holds) $\rightarrow$ hedged carry improves (\textbf{positive, offsetting})
\item CIP basis barely moves $\rightarrow$ manageable
\item \textbf{Net}: forces partially cancel. Self-stabilising.

\end{itemize}
\textit{Stress} (say, a 20\% crash --- March 2020 or a severe global risk-off event):
\begin{itemize}
\item USD appreciates violently (dollar funding squeeze overwhelms all other forces) $\rightarrow$ huge mark-to-market losses, hedging cost surges (\textbf{negative})
\item Yield curve bull-flattens/inverts (flight to quality crushes long-end yields) $\rightarrow$ hedged carry disappears (\textbf{negative, reinforcing} not offsetting)
\item CIP basis blows out $\rightarrow$ hedging prohibitive (\textbf{negative, reinforcing})
\item \textbf{Net}: ALL three forces push in the same direction. Pension fund stops rolling hedges, sells US bonds. Every hedged foreign investor faces the same forces simultaneously $\rightarrow$ \textbf{coordinated sell-off} in US bonds $\rightarrow$ feeds back to equity declines, further USD appreciation, further CIP blowout. This is the positive feedback loop that only activates in the tails.

\end{itemize}
The quantile regression test captures exactly this: the interaction term $\beta_4(\tau) \cdot (\text{Equity}_t \times \text{VIX}_t)$ should be near zero at $\tau = 0.50$ (offsetting forces cancel at the median) but large at $\tau = 0.05$ (all forces reinforce in the tail).


\bigskip\noindent\rule{\textwidth}{0.4pt}\bigskip


\subsubsection{Layer 1: Quantile Regression --- Offsetting Forces Test}

\textbf{Purpose}: Test whether the interaction between equity shocks and FX volatility has a different effect on CIP basis at the median vs. the tails. If ''offsetting forces'' break down in the tails, the interaction coefficient should be near zero at $\tau = 0.50$ but large at $\tau = 0.05$ and $\tau = 0.95$.

\textbf{Unit of observation}: Currency pair \$c\$ $\times$ month \$t\$.

\textbf{Input variables}:


{\small
\begin{longtable}{|p{0.15\textwidth}|p{0.15\textwidth}|p{0.15\textwidth}|p{0.15\textwidth}|p{0.15\textwidth}|p{0.15\textwidth}|}
\hline
\textbf{Variable} & \textbf{Symbol} & \textbf{Definition} & \textbf{Source} & \textbf{Frequency} & \textbf{Level} \\
\hline
\endhead
CIP basis & $\text{CIP}_{c,t}$ & Cross-currency basis swap spread (3M tenor) for currency \$c\$ vs USD & Bloomberg (3M XCCY basis) & Daily $\rightarrow$ Monthly & Currency pair \\
\hline
US equity return & $\text{Equity}_t$ & S\&P 500 monthly log-return & Bloomberg & Daily $\rightarrow$ Monthly & US \\
\hline
Slope differential & $\text{SlopeDiff}_{c,t}$ & (10Y-2Y yield in country \$c\$) minus (10Y-2Y yield in US) & Refinitiv / Bloomberg & Daily $\rightarrow$ Monthly & Currency pair \\
\hline
VIX & $\text{VIX}_t$ & CBOE Volatility Index (monthly average) & FRED (VIXCLS) & Daily $\rightarrow$ Monthly & Global \\
\hline
Interaction & $\text{Equity}_t \times \text{VIX}_t$ & Tests whether equity impact on CIP is amplified during stress & Constructed & Monthly & Global \\
\hline
\end{longtable}
}

\textbf{Currency pairs}: EUR/USD, GBP/USD, JPY/USD, CHF/USD, AUD/USD (G10 majors) + KRW/USD, MXN/USD, BRL/USD, ZAR/USD (key EME pairs with liquid basis data).

\textbf{Model specification}:

\[Q_{\tau}(\text{CIP}_{c,t}) = \alpha_c(\tau) + \beta_1(\tau) \cdot \text{Equity}_t + \beta_2(\tau) \cdot \text{SlopeDiff}_{c,t} + \beta_3(\tau) \cdot \text{VIX}_t + \beta_4(\tau) \cdot (\text{Equity}_t \times \text{VIX}_t) + \varepsilon_{c,t}\]

\textbf{Key parameter}: $\beta_4(\tau)$ --- the interaction between equity shocks and VIX.
\begin{itemize}
\item At $\tau = 0.50$: expect $\beta_4 \approx 0$ (offsetting forces cancel).
\item At $\tau = 0.05$ (left tail) and $\tau = 0.95$ (right tail): expect $|\beta_4|$ large and significant (forces reinforce during extreme events).

\end{itemize}
\textbf{Estimation details}:
\begin{itemize}
\item \texttt{quantile\_cip\_regression()} in \texttt{src/analysis/fx\_hedging\_contagion.py}.
\item \texttt{statsmodels.QuantReg} estimated at $\tau \in \{0.05, 0.10, 0.25, 0.50, 0.75, 0.90, 0.95\}$.
\item Currency fixed effects ($\alpha_c$) included.
\item \texttt{asymmetry\_test()} performs a formal Wald-type test: $H_0: \beta_4(0.05) = \beta_4(0.50)$, using bootstrap standard errors for interquantile inference.

\end{itemize}
\textbf{Output}:
\begin{enumerate}
\item Coefficient table: $\beta_1(\tau), \ldots, \beta_4(\tau)$ across all quantiles with standard errors and p-values.
\item Key comparison: $\beta_4(0.05)$ vs. $\beta_4(0.50)$ --- the ''offsetting forces breakdown'' test.
\item Wald test p-value for interquantile difference.
\item Plot: $\beta_4(\tau)$ across quantiles with confidence bands (should be flat near zero at the median but diverge at the tails).

\end{enumerate}
\textbf{$\rightarrow$ Link to Layer 2}: Layer 1 shows the mechanism exists. Layer 2 quantifies its \textit{systemic risk implication} --- how much does FX hedging stress spill over to bond flows?


\bigskip\noindent\rule{\textwidth}{0.4pt}\bigskip


\subsubsection{Layer 2: Delta-CoVaR (Adrian \& Brunnermeier 2016)}

\textbf{Purpose}: Measure the systemic risk contribution of FX hedging stress to cross-border bond flows. A large negative $\Delta$-CoVaR means that when the CIP basis is in its stressed state, the left tail of bond flows is significantly worse than when the CIP basis is at its median.

\textbf{Unit of observation}: Currency pair \$c\$ $\times$ month \$t\$.

\textbf{Input variables}:


{\small
\begin{longtable}{|p{0.15\textwidth}|p{0.15\textwidth}|p{0.15\textwidth}|p{0.15\textwidth}|p{0.15\textwidth}|p{0.15\textwidth}|}
\hline
\textbf{Variable} & \textbf{Symbol} & \textbf{Definition} & \textbf{Source} & \textbf{Frequency} & \textbf{Level} \\
\hline
\endhead
Bond flows & $F_{c,t}^{Bond}$ & Net cross-border bond fund flows to country \$c\$ & EPFR Global & Monthly & Country \\
\hline
CIP basis & $\text{CIP}_{c,t}$ & Same as Layer 1 & Bloomberg & Monthly & Currency pair \\
\hline
\end{longtable}
}

\textbf{Model specification} --- Two-step quantile regression:

\textbf{Step 1}: Estimate the 5th percentile of the CIP basis:
\[Q_{0.05}(\text{CIP}_{c,t}) = \alpha + \gamma' Z_t\]

This defines the ''stressed'' state: $\text{CIP}_{c,t} = Q_{0.05}(\text{CIP})$.

\textbf{Step 2}: Estimate the 5th percentile of bond flows conditional on the CIP basis:
\[Q_{0.05}(F_{c,t}^{Bond} \mid \text{CIP}_{c,t}) = \alpha_0 + \alpha_1 \cdot \text{CIP}_{c,t}\]

\textbf{Delta-CoVaR} is then:

\[\Delta\text{-CoVaR}_c = Q_{0.05}(F_c^{Bond} \mid \text{CIP}_c = \text{stressed}) - Q_{0.05}(F_c^{Bond} \mid \text{CIP}_c = \text{median})\]

A large negative $\Delta$-CoVaR means FX hedging stress \textit{significantly worsens} the tail risk of bond flows.

\textbf{Estimation details}:
\begin{itemize}
\item \texttt{covar\_fx\_hedging()} in \texttt{src/analysis/fx\_hedging\_contagion.py}. Uses \texttt{statsmodels.QuantReg} at $\tau = 0.05$.
\item \texttt{delta\_covar\_all\_currencies()} computes for all 5+ currency pairs.
\item Confidence intervals via bootstrap (1000 replications).

\end{itemize}
\textbf{Output}:
\begin{enumerate}
\item $\Delta$-CoVaR for each currency pair --- table showing which currencies have the largest systemic risk contribution from FX hedging.
\item Bar chart: $\Delta$-CoVaR ranked by magnitude.
\item Interpretation: e.g., ''When the EUR/USD CIP basis moves from its median (-5 bps) to its 5th percentile (-60 bps), the 5th percentile of European bond flows drops by an additional \$2.5B/month.''

\end{enumerate}
\textbf{$\rightarrow$ Link to Layer 3}: Layer 2 shows the \textit{magnitude} of systemic risk. Layer 3 traces the full \textit{transmission chain} from equities through FX and CIP to bond flows, testing whether the chain is active only during stress.


\bigskip\noindent\rule{\textwidth}{0.4pt}\bigskip


\subsubsection{Layer 3: Quantile Connectedness (Ando et al. 2022)}

\textbf{Purpose}: Estimate the full equity $\rightarrow$ FX $\rightarrow$ CIP $\rightarrow$ bond flows transmission chain as a system, testing whether it is dormant at the median but active during stress.

\textbf{Unit of observation}: System of 4 variables $\times$ month \$t\$ (for each currency pair \$c\$).

\textbf{Input variables}:


{\small
\begin{longtable}{|p{0.15\textwidth}|p{0.15\textwidth}|p{0.15\textwidth}|p{0.15\textwidth}|p{0.15\textwidth}|p{0.15\textwidth}|}
\hline
\textbf{Variable} & \textbf{Symbol} & \textbf{Definition} & \textbf{Source} & \textbf{Frequency} & \textbf{Level} \\
\hline
\endhead
US equity return & $r_t^{EQ}$ & S\&P 500 monthly log-return & Bloomberg & Monthly & US \\
\hline
FX return & $r_{c,t}^{FX}$ & Monthly change in spot exchange rate for currency pair \$c\$ & Bloomberg & Monthly & Currency pair \\
\hline
CIP basis & $\text{CIP}_{c,t}$ & Same as Layer 1 & Bloomberg & Monthly & Currency pair \\
\hline
Bond flows & $F_{c,t}^{Bond}$ & Same as Layer 2 & EPFR & Monthly & Country \\
\hline
\end{longtable}
}

\textbf{System vector}: $Y_t = (r_t^{EQ}, r_{c,t}^{FX}, \text{CIP}_{c,t}, F_{c,t}^{Bond})'$ --- a \textbf{4-variable Quantile VAR} representing the transmission chain for a single currency pair.

\textbf{Model specification} --- Same QVAR framework as P2 Module 3:

\[Q_{\tau}(y_{i,t} \mid Y_{t-1}, \ldots, Y_{t-p}) = c_i(\tau) + \sum_{j=1}^{4} \sum_{l=1}^{p} \beta_{ij}^{(l)}(\tau) \, y_{j,t-l}\]

\textbf{Estimation details}:
\begin{itemize}
\item \texttt{build\_connectedness\_system()} constructs the 4-variable system from the synthetic or real data.
\item \texttt{fx\_hedging\_quantile\_connectedness()} estimates the QVAR and computes GFEVD at each quantile.
\item \texttt{tail\_connectedness\_comparison()} compares total connectedness across $\tau \in \{0.05, 0.50, 0.95\}$.
\item \texttt{cross\_currency\_connectedness()} repeats for all currency pairs.
\item Lags: \$p = 4\$; horizon: \$h = 10\$.

\end{itemize}
\textbf{Key prediction}: The transmission chain equity $\rightarrow$ FX $\rightarrow$ CIP $\rightarrow$ bond flows is:
\begin{itemize}
\item \textbf{Dormant at $\tau = 0.50$}: low total connectedness, weak off-diagonal GFEVD entries.
\item \textbf{Active at $\tau = 0.05$}: high total connectedness, strong sequential spillovers from equity to FX to CIP to bond flows.

\end{itemize}
\textbf{Output}:
\begin{enumerate}
\item Total connectedness at each quantile: \$C(0.05)\$, \$C(0.50)\$, \$C(0.95)\$ --- for each currency pair.
\item GFEVD heatmaps at $\tau = 0.05$ vs. \$0.50\$: the $4 \times 4$ matrix showing pairwise spillovers.
\item Directional analysis: equity is the dominant net transmitter in the left tail; bond flows are the dominant net receiver.
\item Cross-currency summary table: which currency pairs show the strongest tail amplification.

\end{enumerate}

\bigskip\noindent\rule{\textwidth}{0.4pt}\bigskip


\subsubsection{NBFI Penetration $\times$ FX Hedging (Panel Regression)}

\textbf{Purpose}: Test whether countries with higher NBFI penetration suffer stronger contagion through the FX hedging channel.

\textbf{Unit of observation}: Country \$c\$ $\times$ month \$t\$.

\textbf{Input variables}:


{\small
\begin{longtable}{|p{0.15\textwidth}|p{0.15\textwidth}|p{0.15\textwidth}|p{0.15\textwidth}|p{0.15\textwidth}|p{0.15\textwidth}|}
\hline
\textbf{Variable} & \textbf{Symbol} & \textbf{Definition} & \textbf{Source} & \textbf{Frequency} & \textbf{Level} \\
\hline
\endhead
Bond flows & $F_{c,t}^{Bond}$ & Net cross-border bond fund flows & EPFR & Monthly & Country \\
\hline
CIP basis & $\text{CIP}_{c,t}$ & Cross-currency basis swap spread & Bloomberg & Monthly & Currency pair \\
\hline
NBFI penetration & $\text{NBFI}_{c,t}$ & NBFI assets as \% of GDP & FSB & Annual (interpolated) & Country \\
\hline
Interaction & $\text{CIP}_{c,t} \times \text{NBFI}_{c,t}$ & Tests amplification & Constructed & Monthly & Country \\
\hline
Controls & $X_{c,t}$ & VIX, local equity return, yield slope differential & Various & Monthly & Mixed \\
\hline
\end{longtable}
}

\textbf{Model specification --- OLS panel}:

\[F_{c,t}^{Bond} = \alpha_c + \beta_1 \cdot \text{CIP}_{c,t} + \beta_2 \cdot \text{NBFI}_{c,t} + \beta_3 \cdot (\text{CIP}_{c,t} \times \text{NBFI}_{c,t}) + \gamma' X_{c,t} + \varepsilon_{c,t}\]

\textbf{Key parameter}: $\beta_3 < 0$. Countries with higher NBFI penetration experience \textit{larger} bond flow reversals when CIP basis widens (hedging costs spike).

\textbf{NBFI counterparty decomposition}: As in Project 3, the aggregate $\text{NBFI}_{c,t}$ should be decomposed by counterparty type. In this context, the most relevant decomposition is by \textbf{FX hedging activity}: countries where NBFIs rely heavily on cross-border FX-hedged bond holdings (foreign-borrowing model) should show the strongest amplification, while countries where NBFIs are primarily domestic credit intermediaries should be less affected by the CIP channel. The from-whom-to-whom financial accounts provide the basis for this decomposition.

\textbf{IV specification}: CIP basis instrumented with central bank swap line announcements (exogenous supply-side shocks to FX funding).

\textbf{Estimation details}:
\begin{itemize}
\item \texttt{nbfi\_fx\_hedging\_panel\_regression()} and \texttt{iv\_fx\_hedging\_regression()} in \texttt{src/analysis/fx\_hedging\_contagion.py}.
\item Country FE, clustered standard errors.

\end{itemize}
\textbf{Output}: Coefficient table (OLS + IV); scatter plot split by high/low NBFI.


\bigskip\noindent\rule{\textwidth}{0.4pt}\bigskip


\subsubsection{Cascade Simulation: FX Hedging Feedback Loop}

\textbf{Purpose}: Simulate the full feedback loop: equity shock $\rightarrow$ FX pass-through $\rightarrow$ CIP widening $\rightarrow$ bond flow reduction $\rightarrow$ equity feedback $\rightarrow$ VIX spike $\rightarrow$ repeat.

\textbf{Unit of observation}: Simulated rounds (no agent-based --- this is a reduced-form cascade).

\textbf{Input variables (calibration)}:


{\small
\begin{longtable}{|p{0.15\textwidth}|p{0.15\textwidth}|p{0.33\textwidth}|p{0.22\textwidth}|}
\hline
\textbf{Parameter} & \textbf{Symbol} & \textbf{Value} & \textbf{Rationale} \\
\hline
\endhead
Initial equity shock & --- & -5\% & Scenario \\
\hline
FX pass-through & $\phi$ & 0.30 & 30\% of equity shock passes to FX (Rey et al.) \\
\hline
CIP sensitivity & --- & 5 bps per 1\% FX move & Calibrated from BIS data \\
\hline
Bond flow sensitivity & --- & -2.0 (\$B per 10 bps CIP) & EPFR regression coefficients \\
\hline
Equity feedback & --- & 0.15 (15\% of bond flow shock feeds back to equity) & Portfolio rebalancing effect \\
\hline
Dampening & --- & $0.7^{(r-1)}$ per round & Geometric decay (markets adjust) \\
\hline
Max rounds & --- & 20 & Until convergence \\
\hline
\end{longtable}
}

\textbf{Simulation mechanism}:
\begin{enumerate}
\item \textbf{Round 1}: Equity falls -5\%.
\item \textbf{FX}: USD appreciates by $0.30 \times 5\% = 1.5\%$.
\item \textbf{CIP}: Basis widens by $5 \times 1.5 = 7.5$ bps.
\item \textbf{Bond flows}: Reverse by $(-2.0) \times 7.5/10 = -\$1.5B.
\item \textbf{Feedback}: Equity drops additional $0.15 \times 1.5/100 = 0.225\%$.
\item \textbf{VIX}: Spikes proportional to cumulative equity loss.
\item \textbf{Dampening}: Each subsequent round is 70\% of the previous.
\item \textbf{Repeat} until changes $< 0.01\%$.

\end{enumerate}
\textbf{Estimation details}: \texttt{simulate\_fx\_hedging\_cascade()} in \texttt{src/analysis/fx\_hedging\_contagion.py}.

\textbf{Output}:
\begin{enumerate}
\item Round-by-round dynamics: equity return, FX move, CIP change, bond flow, VIX change, cumulative loss.
\item Total cascade multiplier: cumulative loss / initial shock. If \$> 1\$: the FX hedging channel amplifies.
\item Convergence chart: how quickly the feedback loop dampens.

\end{enumerate}

\bigskip\noindent\rule{\textwidth}{0.4pt}\bigskip


\subsubsection{Module Linkage Map (Project 4)}


\begin{adjustbox}{max width=\textwidth}
\begin{minipage}{1.5\textwidth}
\begin{verbatim}
+--------------------------------------+
| Layer 1: Quantile Regression         |
| CIP = f(Equity, Slope, VIX, EqxVIX) |
| Does interaction break down in tails?|
+--------------+-----------------------+
               | ''Yes --- offsetting forces fail in tails''
               v
+--------------------------------------+
| Layer 2: Delta-CoVaR                 |
| Systemic risk: CIP stress → bond    |
| flow left tail                       |
| ''How bad is it?''                     |
+--------------+-----------------------+
               | ''CIP stress significantly worsens bond flow tails''
               v
+--------------------------------------+
| Layer 3: Quantile Connectedness      |
| Full chain: Equity → FX → CIP →     |
| Bond flows (QVAR on 4-var system)    |
| ''Chain dormant at median, active     |
|  at tau = 0.05''                        |
+----------+-----------+---------------+
           |           |
           v           v
+----------------+  +---------------------+
| NBFI Panel     |  | Cascade Simulation  |
| CIP x NBFI    |  | Equity → FX → CIP → |
| amplification  |  | Bonds → Feedback    |
| (OLS + IV)     |  | Multi-round         |
+----------------+  +---------------------+
\end{verbatim}
\end{minipage}
\end{adjustbox}


\subsection{Expected Contributions}

\begin{enumerate}
\item First complete model of the equity-FX-hedging-bonds contagion chain, combining Rey et al. (2024) and Nenova et al. (2025).
\item Discovery of the 'offsetting forces' mechanism: partial cancellation at the median, reinforcement in the tails.
\item First application of quantile connectedness to the FX hedging transmission chain.
\item Evidence that NBFI FX hedging maturity mismatch (3-month swaps hedging 10-year bonds) creates systemic rollover risk.

\end{enumerate}

\bigskip\noindent\rule{\textwidth}{0.4pt}\bigskip


\section{Common Methodological Toolkit}

All four projects share a core econometric toolkit:


{\small
\begin{longtable}{|p{0.22\textwidth}|p{0.33\textwidth}|p{0.33\textwidth}|}
\hline
\textbf{Method} & \textbf{What It Captures} & \textbf{Used In} \\
\hline
\endhead
Quantile Regression & Asymmetric effects across the distribution & P1, P2, P3, P4 \\
\hline
Quantile VAR (QVAR) & Tail-specific shock propagation & P2, P3, P4 \\
\hline
Quantile Connectedness & Tail spillovers between variables/countries & P1, P2, P3, P4 \\
\hline
Delta-CoVaR & Systemic risk contribution of specific channels & P1, P4 \\
\hline
Location-Scale GaR & Mean + variance effects on future growth & P2, P3 \\
\hline
DCC-GARCH & Time-varying correlations across sectors & P2 \\
\hline
Diebold-Yilmaz GFEVD & Directional connectedness (mean) & P2, P3 \\
\hline
Network / spectral & Centrality, HLI, contagion matrices (credit + funding) & P1, P3 \\
\hline
Agent-based simulation & Cascade / fire-sale propagation & P1, P3, P4 \\
\hline
Panel FE + IV (2SLS) & Causal identification with MP shocks & P3, P4 \\
\hline
Quantile factor model & Global Tail Cycle extraction (Ando \& Bai 2020) & P2 (used in P3, P4) \\
\hline
From-whom-to-whom accounts & NBFI counterparty decomposition & P1, P3, P4 \\
\hline
\end{longtable}
}


\subsubsection{Cross-Cutting Analytical Themes}

Three themes recur across all four projects and deserve explicit recognition:


\paragraph{1. State-Dependence and Distance to Constraints}

The propagation of any shock through the NBFI system depends on two things simultaneously:
\begin{itemize}
\item \textbf{Structural integration}: How embedded is NBFI activity in the financial system (size, leverage, interconnectedness)? This is relatively slow-moving.
\item \textbf{Distance to constraints}: How close are intermediaries to their VaR/margin/redemption constraints at the moment the shock hits? This is fast-moving and state-dependent.

\end{itemize}
A moderate monetary policy tightening when hedge funds have ample excess margin is very different from the same shock when funds are already near VaR limits. Same structural integration, dramatically different propagation. The QVAR captures this implicitly ($\beta_{ij}(\tau)$ varies with $\tau$), but it can be made explicit via threshold specifications or by interacting NBFI variables with distance-to-constraint measures.

The dangerous combination is: high $\lambda_1(W_t)$ (tight network) \textbf{plus} low slack (near constraints). Neither alone is sufficient for systemic crisis.


\paragraph{2. NBFI Counterparty Decomposition}

Aggregate ''NBFI penetration'' (assets / GDP) conflates activities with fundamentally different systemic implications. The key decomposition is by counterparty:


{\small
\begin{longtable}{|p{0.22\textwidth}|p{0.33\textwidth}|p{0.33\textwidth}|}
\hline
\textbf{Domestic banks} & \textbf{Foreign counterparties} & \textbf{Domestic private sector} \\
\hline
\endhead
\textbf{NBFI lends to} & Funding source for banks & Transmits shocks outward \\
\hline
\textbf{NBFI borrows from} & Bank credit risk exposure & Imports foreign funding shocks \\
\hline
\end{longtable}
}

This decomposition matters for every panel regression in the programme. Data: from-whom-to-whom financial accounts (Fed Z.1, ECB, OECD).


\paragraph{3. Global Tail Cycle}

The standard Global Financial Cycle (Rey 2013) captures co-movement at the mean. A \textbf{Global Tail Cycle} captures co-movement in the left tail --- the common factor in crash risk. These are conceptually distinct objects: markets can look diversified on average but crash together. The GTC is the risk factor that sets in motion coordinated VaR breaches and feedback loops across jurisdictions. \textbf{Constructed in Project 2, Module 3} (QVAR) via the Ando \& Bai (2020) quantile factor model at $\tau = 0.05$, or alternatively as the first principal component of the QVAR tail spillover matrix. Used as an input in Project 3 (Module 3) and Project 4 (Module 5).


\bigskip\noindent\rule{\textwidth}{0.4pt}\bigskip


\section{Cross-Project Linkages}

The four projects are complementary. Each addresses a distinct dimension of NBFI
systemic risk, but findings feed into each other:

\begin{itemize}
\item \textbf{P1 $\rightarrow$ P2}: Implied leverage estimates serve as inputs to P2's VaR-constraint
channel. Higher hidden leverage = tighter VaR constraints = amplified MP transmission.

\item \textbf{P2 $\rightarrow$ P3}: P2 identifies which NBFI sub-sectors amplify domestic MP transmission;
P3 traces how this amplification propagates across borders via dollar funding
and portfolio flow channels.

\item \textbf{P3 $\rightarrow$ P4}: P3's dollar funding channel (FX swap basis) is the same mechanism
P4 models in detail. P4 adds the equity-FX link (Rey et al. 2024) and the yield
curve offset (Nenova et al. 2025).

\item \textbf{P4 $\rightarrow$ P1}: FX derivatives are a major source of hidden leverage (P1), and P4
shows how hedging activity in these instruments creates contagion. Together they
demonstrate that derivatives are simultaneously a source of hidden leverage AND
a contagion channel.

\end{itemize}

\begin{quote}
\itshape
\textbf{Common finding across all four projects}: NBFI risks are fundamentally \\
nonlinear. Standard mean-based methods miss the tail amplification, asymmetric \\
contagion, and crisis-activated transmission that quantile methods reveal. \\
\end{quote}

\begin{itemize}
\item \textbf{P3 $\leftrightarrow$ Global Tail Cycle}: P3's GFC factor extraction is extended to a
Global \textit{Tail} Cycle --- a common factor in crash risk rather than average
co-movement. NBFI penetration amplifies GTC transmission more than GFC
transmission, confirming that NBFI systemic risk is fundamentally a tail
phenomenon.

\item \textbf{Cross-cutting: NBFI counterparty decomposition}: Across P1, P3, and P4,
aggregate NBFI measures are decomposed by counterparty type (bank-lending,
bank-borrowing, foreign, private sector). This reveals which \textit{type} of NBFI
activity drives amplification --- not just how \textit{much} NBFI there is.

\item \textbf{Cross-cutting: State-dependence}: Across P1--P4, shock amplification
depends not only on structural NBFI integration but on distance to
constraints. The HLI (P1) $\times$ constraint slack interaction captures
the joint condition for systemic crisis.
\end{itemize}

\end{document}
